\documentclass[final,5p,times,twocolumn,authoryear]{elsarticle}

% \documentclass[final,5p,times,twocolumn,authoryear]{autart}     % TODO CRITICAL

% TODO: You need to change to \documentclass[twocolumn]{autart}

\include{Config/config_automatica}
\include{Config/config_des}

\journal{Automatica}

\begin{document}

\begin{frontmatter}

\title{Synthesis of Cognitive Non-Blocking Supervisors for Partially Observed Discrete-Event Systems}

\author[inst1]{Yiwei Zheng}
\ead{ywzheng1201@sjtu.edu.com}
\author[inst1]{Bohan Cui}
\ead{bohan\_cui@sjtu.edu.cn}
\author[inst1]{Xiang Yin\corref{cor1}}
\ead{yinxiang@sjtu.edu.cn}

\cortext[cor1]{Corresponding author.}
\address[inst1]{TODO Department of Example Engineering, Example University, City, Country}

\begin{abstract}
This article investigates the problem of synthesizing cognitive non-blocking supervisors for partially observed discrete event systems (DESs).
The notion of cognitive non-blockingness requires that for each state of the DES, the system can eventually certify its arrival at a target state based solely on partial observations.
The synthesis problem is therefore to construct a supervisor under which the closed-loop system satisfies this observation-based progress requirement.
%
We show that the non-blocking all-inclusive controller (AIC) proposed in \cite{yin2016synthesis} is insufficient for this purpose, and introduce a new structure, named the bipartite cognitive controller (BCC), to capture the microscopic evolution within the AIC.
We prove that the event sequences generated by a BCC can be directly projected as the language of the controlled DES, thereby establishing the equivalence between the BCC-induced language and the controlled DES language.
Moreover, we show that cognitive non-blockingness can be guaranteed by requiring the decision states to be cognitively live and the observation states to be potentially feasible.
The resulting BCC can then be projected back to a supervisor for cognitive non-blocking synthesis of the original DES.
\end{abstract}

\begin{keyword}
Discrete-event systems (DES) \sep partial observation \sep non-blocking, supervisory control
\end{keyword}

\end{frontmatter}

\section{Introduction}
\label{sec:intro}
As modern computation and networking technologies continue to advance, discrete event systems (DESs) have become increasingly important in various aspects of modern life and industrial production \cite{yin2024formal, fokkink2024offline}.
Owing to their capability to model complex event-driven behaviors, DESs have been widely applied in areas such as automated manufacturing systems, communication networks, transportation systems, and robotic coordination.
%
Non-blockingness is one of the most fundamental properties of DESs, since it guarantees that the system can always complete its prescribed tasks without being trapped before task completion \cite{cassandras2021introduction}.
In general, a DES is said to be non-blocking if every reachable state can still evolve to a marked (or target) state through some feasible event sequence.
This ensures that the system can continuously make progress toward task completion rather than entering an irreversible blocking state, for example, a deadlock or a livelock state.

%
% TODO
% 要不要把enforce这个词引入?
The supervisory control framework introduced by Ramadge and Wonham \cite{ramadge1987supervisory} provides a systematic approach to enforcing desired behaviors of DESs by enabling or disabling controllable events.
In the classical non-blocking synthesis problem, the objective is to construct a supervisor such that every behavior permitted by the closed-loop system can be extended to a marked behavior.
%
However, under partial observation, classical non-blockingness does not guarantee observational certifiability of reaching marked states. Even if a marked state is reachable in the underlying system, the supervisor may be unable to establish this based on partial observations.
%
Let us consider the motivating example of a robotic system deployed in a partially observable environment, which is illustrated in Fig.~\ref{fig:dfa_grid_env} (a). 
The robot has full knowledge of the environment model, while its exact location remains uncertain due to insufficient sensing and limited localization accuracy.
Suppose that the robot is allowed to transmit data only in safe regions, which are regarded as the target regions in this task. 
If a safe target location is observationally indistinguishable from an unsafe location, then physical reachability of a safe region is not sufficient for safe data transmission.
Instead, safety must be certifiable from observations.
Therefore, before transmitting data, the robot must ensure that all states consistent with its observation belong to the set of safe target locations.

This mismatch between non-blockingness and observational certifiability motivates us to introduce the notion of \emph{cognitive non-blockingness}, an observation-based progress property requiring that task completion be not only reachable but also certifiable under partial observation.
%
We then study the synthesis of partial-observation supervisors that enforce cognitive non-blockingness on the closed-loop system.
%
Built upon the supervisory control framework, this article further studies the synthesis of partial-observation supervisors that enforce cognitive non-blockingness on the closed-loop system.
Different from classical non-blocking synthesis, the proposed problem must account for the evolution of information states, since task completion has to be certified from observations rather than merely reached by the actual state.

%
% TODO 争取给放到第二页?
\begin{figure}[t!]
    \centering
    \includegraphics[width=\linewidth]{Figs/Grid_Environment_DFA.pdf}
    \caption{Deterministic finite state automaton $G$ constructed from the motivating example of the robotic system deployed in a partially observable environment from Section~\ref{sec:intro}, where the set of marked states $Q_m = \{ q_3, q_5, q_6, q_{18} \}$, the set of observable events $\Sigma_o = \{ o_1, o_2, o_3 \}$, and the set of controllable events $\Sigma_c = \{ c_1, c_2, o_3 \}$.}
    % \textcolor{red}{not good, the feasible transitions in (a) are not fully captured by (b), e.g., 19 $\to$ 13, 19 $\to$ 15..., maybe show them separately or motivate the paper only use words}
    %
    % Thanks for your comment
    %
    \label{fig:dfa_grid_env}
\end{figure}

The main contributions of this article are summarized as follows.
%
\begin{enumerate}[(i)]
    \item We propose a new observation-based property, named \emph{cognitive non-blocking}, for partially observed DESs.
    Compared with classical non-blockingness studied in  \cite{yin2016synthesis, yin2018synthesis}, the proposed notion strengthens the standard reachability-based requirement by incorporating observational certification of task completion.
    Specifically, it requires not only that marked states are eventually reachable, but also that the supervisor can certify, based solely on partial observations, that the current state belongs to the marked set.

    \item We develop a new structural representation, called \emph{bipartite cognitive controller} (BCC), to characterize both the evolution of macrostate-level and the plant-state-level transitions hidden within each unobservable reach. 
    Existing bipartite-controller-based approaches \cite{yin2016uniform,zheng2026enforcement} describe closed-loop evolution in terms of control decisions and observations. In contrast, the proposed BCC further refines this representation at the plant-state level, thereby enabling the synthesis of cognitive non-blocking supervisors.

    \item We further develop a synthesis method for cognitive non-blocking supervisors based on the BCC structure. 
    This result strengthens the framework of \cite{yin2016synthesis} by iteratively pruning observation states that are not potentially feasible and decision states that are not cognitively live, while preserving structural consistency, so that the closed-loop system guarantees cognitive non-blockingness for every reachable state.
    
\end{enumerate}

\subsection{Related works}

% 1 partially observable supervisor for the IS-based property
The supervisory control of DESs under partial observation has been extensively studied in the context of information-state-based (IS-based) properties such as safety \cite{ma2022resilient}, diagnosability \cite{li2025supervisory, hu2023asynchronous}, opacity \cite{lafortune2018history, zheng2026enforcement},
detectability \cite{shu2013enforcing}, anonymity and attractability, etc. \cite{yin2024formal, basile2022supervisory}.
These properties require the synthesis of a supervisor that ensures that all executions of the closed-loop system satisfy the given property.
%
In \cite{yin2016uniform}, a uniform approach is developed to synthesize a bipartite controller that embeds all reachable information states and all admissible supervisory control strategies, which can be further used to synthesize the maximally permissive supervisor for specific properties.
In \cite{xie2023optimal}, this bipartite controller is extended to solve the problem of enforcing opacity for both qualitative (or infinite-step opacity) and quantitative cases.
In \cite{zheng2026enforcement}, an extended bipartite controller is developed for interval weighted DESs for the enforcement of weighted opacity, for which the enabling interval and the sequences of control patterns determine the reachability of the system.
%
However, most existing IS-based properties focus on information-state consistency, where system behavior can be characterized without explicitly reasoning about future progress. 
In contrast, non-blockingness is inherently a progress property that depends on the existence of continuations to marked states and therefore requires reasoning beyond current information states. 
As a result, directly pruning inconsistent information states using classical bipartite controllers is insufficient, which motivates the development of new observation-based progress properties.


% 2 non-blocking problems
Non-blockingness is a fundamental progress requirement in DESs, requiring that every generated behavior can be extended to a marked behavior \cite{chen1991dealing}.
%
The verification of non-blockingness has been extensively investigated for different DES frameworks, including modular DESs \cite{tang2024compositional}, adapted systems \cite{thuijsman2024transformational}, and compositional analysis of extended finite-state machines \cite{mohajerani2016framework}.
%
In \cite{tomas2018complexity}, the computational complexity of verifying non-blockingness in modular DESs is rigorously analyzed, showing that the problem is PSPACE-complete in general, while special cases such as DFA with one shared event or a bounded number of DFAs are NP- or NL-complete.
%
In \cite{pilbrow2015algorithm}, the authors develop a set of compositional abstraction rules, including blocked events, to efficiently verify non-blockingness in modular DESs.
%
In \cite{tang2024compositional}, a framework is developed for compositional non-blocking verification of finite automata with prioritized events, introducing conflict-preserving abstractions and shaped quotient automata to efficiently ensure non-blockingness in modular DESs.
%
However, most of these works on non-blocking analysis focus on verification rather than synthesis. In many practical scenarios, it is necessary to synthesize a supervisor to actively control the system such that its output behaviors are guaranteed to satisfy the non-blocking property \cite{chen1991dealing, malik2013compositional}. 
This motivates the development of synthesis methods that ensure non-blockingness under various system architectures and observation constraints, bridging the gap between theoretical verification results and practical supervisory control design.


% 3 A + B 
In recent research, the combination of non-blockingness and partial-observation supervisory control has also been widely investigated \cite{fokkink2024offline, zhang2024quantitatively}.
%
In \cite{yin2016synthesis}, the IS-based bipartite controller of \cite{yin2016uniform} is extended to synthesize the safe, non-blocking, and maximally permissive supervisor.
%
In \cite{xiang2024existence}, a game-theoretic framework is developed to synthesize bounded non-blocking supervisors for DESs by transforming the supervisor existence problem into a reachability game over information states and admissible control decisions.
%
In \cite{jan2024supervisory}, modular supervisory control under partial observation is investigated from the perspective of normality, where a global non-blocking and maximally permissive normal supervisor is achieved through the parallel composition of local normal supervisors.
%
In \cite{zhang2025bounded},  a new bounded-time non-blocking supervisory control problem is solved for timed DES, in which each task must be completed within a bounded time counted by the number of ticks.
%
Although significant progress has been made in integrating non-blockingness with partial-observation supervisory control \cite{fokkink2024offline, zhang2024quantitatively}, non-blockingness only guarantees that the system can eventually reach a marked state corresponding to task completion. 
This reveals a fundamental gap between reachability under supervisory control and its certifiability under partial observation, which motivates the introduction of observation-based non-blocking properties that explicitly connect system observations to task progress.


\subsection{Organization}

The remainder of this article is organized as follows.
%
Section~\ref{sec:prelimi} presents the system model and the notions for partial-observation supervisory control.
Section~\ref{sec:prob} introduces the notion of cognitive non-blockingness, gives the problem formulation, and illustrates the challenges of the proposed problem.
Section~\ref{sec:bcc} puts forward the new structure of the  bipartite cognitive controller and establishes its basic properties.
Section~\ref{sec:solu} presents the method for synthesizing supervisors that enforce cognitive non-blockingness.
Section~\ref{sec:special_z_marked} discusses the special case in which the marked states are within the unobservable reach.
Section~\ref{sec:con} concludes the article and gives directions for future research.

\section{Preliminaries}\label{sec:prelimi}

This section introduces the preliminaries for the verification and synthesis of cognitive non-blocking in discrete-event systems (DESs), including system models and supervisory control under partial observation.


\subsection{System model}

We use $\mathbb{N}$, $\mathbb{N}_{\geq 0}$ to denote the sets of natural numbers and nonnegative natural numbers, respectively. 
We use $\varepsilon$ to denote the empty string of an automaton. 
%

We consider a DES  modeled as a \emph{deterministic} finite-state automaton (DFA)  
    $$
        G = (Q, \Sigma, \delta, q_0, Q_m)
    $$
    where
    $Q$ is the finite set of states, 
    $\Sigma$ is the finite set of events, 
    $\delta: Q \times \Sigma \to Q$ is the   transition function,
    $q_0\in Q$ is the initial state, 
    and $Q_m\subseteq Q$ is the set of marked states. 
The transition function is also extended to 
$\delta:Q\times \Sigma^* \to Q$ recursively in the usual manner. 
The language generated by $G$ is defined by 
$\mathcal{L}(G) = \{ s \in \Sigma^* \mid \delta(q_0, s)! \}$, where $!$ means ``is defined'', and the marked language of $G$ is 
$\mathcal{L}_m(G) = \{ s \in \Sigma^* \mid \delta(q_0, s) \in Q_m \}$. 

%
% 在这里无需再讨论what is the sequence of DFAs looks like? 
% 所以rho就可以被提取到后面用了
In the partial observation setting, the event set $\Sigma$ is further  partitioned into two disjoint subsets 
\[
\Sigma = \Sigma_o \ \dot{\cup} \ \Sigma_{uo},
\]
where $\Sigma_o$ is the set of observable events, and $\Sigma_{uo}$ is the set of unobservable events. 
The \emph{natural projection} $P : \Sigma^* \rightarrow \Sigma_{o}^*$ is recursively defined as
\begin{equation}
P(\varepsilon)  = \varepsilon\text{ and }
P(s \sigma) =
\begin{cases}
P(s), & \text{if } \sigma \in \Sigma_{uo} \\
P(s)\sigma, & \text{otherwise}
\end{cases}
\label{eq:proj}
\end{equation} 
The natural projection is also extended to $P:2^{\Sigma^*}\to 2^{\Sigma_o^*}$ by $P(L)=\{P(s)\in \Sigma_o^* \mid s\in L\}$  for any $L\subseteq\Sigma^*$. 
    
%
% The \emph{run} of DFA $G$ is a sequence of states and transitions
% $
% \rho = q_0 \xrightarrow[]{\sigma_0} q_1 \xrightarrow[]{\sigma_1} \cdots q_n \xrightarrow[]{\sigma_n} q_{n+1},
% $
% where $(q_{i}, \sigma_i, q_{i+1}) \in \delta$ for $0 \leq i \leq n, n \in \mathbb{N}_{\geq 0}$.
% The \emph{sequence} corresponding to run $\rho$ is denoted as
% $tr(\rho) = \sigma_0 \, \sigma_1 \, \cdots \, \sigma_n \in \Sigma^*$,
% %
% With a slight abuse of notation, we use $\delta(q, r) = q'$ as the extended partial transition function to denote the reachable state $q'$ from state $q$ and sequence $r$.
% %
% The behavior generated by $G$ is described by 
% $\mathcal{L}(G) = \{ s \in \Sigma^* | \delta(q_0, s)! \}$, 
% where $!$ means ``is defined'', and the marked behavior is 
% $\mathcal{L}_m(G) = \{ s \in \Sigma^* | \delta(q_0, s) \in Q_m \}$. 
% The prefix-closure of a language $\mathcal{L}(G)$ is
% $$
% \overline{\mathcal{L}(G)} = \{ s \in \Sigma^* | \exists t \in \Sigma^* : st \in \mathcal{L}(G) \},
% $$
% and we say that $\mathcal{L}$ is prefix-closed if $\overline{\mathcal{L}} = \mathcal{L}$.

% In DFA $G$, a state $q \in Q$ is said to be \emph{accessible} if there exists a sequence $s \in \mathcal{L}(G)$ such that $\delta(q_0, s) = q$. 
% A state $q \in Q$ is \emph{co-accessible} if there exists a sequence $t \in \mathcal{L}(G)$ such that $\delta(q, t) \in Q_m$.
% %
% A standard progress property in DES is \emph{non-blocking}, requiring that every reachable state can eventually reach a marked state, as formally given in the following definition.
% \begin{definition}
%     Given a DFA $G=(Q,\Sigma,\delta,q_0,Q_m)$, we say that $G$ is
%     \emph{non-blocking} if every accessible state of $G$ is co-accessible to  
%     $Q_m$,
%     \begin{equation} \label{eq:nb_1}
%         \forall s \in \mathcal{L}(G), \exists t \in \Sigma^* \ s.t. \ st \in \mathcal{L}_m(G).
%     \end{equation}
%     \label{def:nb}
% \end{definition}
% %
% Based on Definition \ref{def:nb}, one can easily check that the system is non-blocking if and only if 
% \begin{flalign}
%     \forall q\in Q, \
%     &(\exists s \in \mathcal{L}(G): \delta(q_0,s)=q)
%     \Rightarrow  \notag \\
%     &(\exists t\in\Sigma^*: \delta(q,t)\in Q_m).
% \label{eq:nb_2}
% \end{flalign}
% Intuitively, the notion of non-blockingness means that the system can always continue from any reachable state to complete its task. 
% In other words, no reachable behavior is trapped before reaching a marked state.



\subsection{Supervisory Control under Partial Observations}

In the supervisory control framework, 
the event set $\Sigma$ is further partitioned into two disjoint subsets
\[
\Sigma = \Sigma_c \ \dot{\cup} \ \Sigma_{uc},
\]
where $\Sigma_c$ is the set of controllable events, and $\Sigma_{uc}$ is the set of uncontrollable events. 
The set of \emph{control patterns} is defined by 
\[
\Gamma=\{\gamma\in 2^\Sigma \mid \Sigma_{uc}\subseteq \gamma \}, 
\]
i.e., all uncontrollable events are always enabled. 
Then a partial-observation supervisor $S$ is a mapping 
\[
S : P(\mathcal{L}(G)) \to  \Gamma
\]
that dynamically enables/disables events based on its observation. 

We use $S / G$ to denote controlled system, and the language generated by $S /G$, denoted by $\mathcal{L}(S / G)$, is recursively defined as follows,
\begin{enumerate}[(i)]
    \item $\varepsilon \in \mathcal{L}(S / G)$;
    \item $s \, \sigma \in \mathcal{L}(S / G)$ if and only if
    $s \in \mathcal{L}(S / G)$ $\wedge$ $s \, \sigma \in \mathcal{L}(G)$ $\wedge$ $\sigma \in S(P(s))$.
\end{enumerate}
We define the marked language of the closed-loop system  as $\mathcal{L}_m(S / G) = \mathcal{L}(S / G) \cap \mathcal{L}_m(G)$.




%
It is worth noting that the desired closed-loop property can be achieved by a trivial supervisor that disables controllable events whenever possible, but such a solution is overly conservative and sacrifices behavioral permissiveness. Therefore, we seek supervisors that satisfy the specification while preserving maximal admissible behavior.
%
A supervisor $S$ is said to be maximally permissive within a prescribed class of admissible supervisors if there does not exist another admissible supervisor $S'$ such that
\begin{equation}\label{eq:mps}
    \mathcal{L}(S/G) \not \subset \mathcal{L}(S'/G).
\end{equation}
%
For a sequence $s\in\mathcal{L}(S/G)$, the \emph{immediate state estimate}, namely, the set of states in which the system could be immediately after the last observed event in $P(s) \in P(\mathcal{L}(G))$, is defined as
\begin{equation}\label{eq:est_o}
    \begin{aligned}
    \hat{\mathcal{E}}_S(\alpha) := \{ 
    &q \in Q \mid \exists \omega \in (\Sigma^*\Sigma_o) \cap \mathcal{L}(S / G) : \\ 
    &\delta(q_0, \omega) = q \wedge P(\omega) = \alpha \},
    \end{aligned}
\end{equation}
and the corresponding \emph{state estimate with unobservable tail} is defined as
\begin{flalign}\label{eq:esti:1}
\mathcal{E}_S(s) := \{ &q \in Q \mid \exists t \in \mathcal{L}(S / G) : \notag \\
                       &\delta(q_0, t) = q \wedge P(t) = P(s)\}.
\end{flalign}



\begin{remark}\label{rem:uc_ctrl_pattern}
    % In DFAs, in accordance with the Ramadge-Wonham framework of \textup{\cite{ramadge1987supervisory}}, uncontrollable events are always enabled, regardless of the control patterns. 
    % That is, transitions labeled by uncontrollable events are guaranteed to be accessible regardless of the control policy.
    % %
    % In articles such as \cite{yin2016synthesis} and \cite{yin2018synthesis}, it is required that $\Sigma_{uc} \subseteq \gamma$, $\forall \gamma \in \Gamma$.
    % %
    % To highlight the controllability aspect of the system, we omit uncontrollable events from the control specifications $\gamma$ and $\Gamma$.
    % %
    % Their accessibility is treated as implicitly enabled in the remainder of this article.
    
    In the standard Ramadge-Wonham framework \cite{ramadge1987supervisory}, control patterns are often defined as subsets of $\Sigma$, and the uncontrollable events are always enabled by the supervisor regardless of the specific control pattern, i.e., $\Sigma_{uc}\subseteq \gamma$ for all $\gamma\in\Gamma$, as in \cite{yin2016synthesis,yin2018synthesis}.
    In this article, we adopt an equivalent convention by defining $\Gamma=2^{\Sigma_c}$, so that each $\gamma\in\Gamma$ records only the enabled controllable events.
    Uncontrollable events are treated as implicitly enabled throughout; equivalently, the actual enabled event set associated with $\gamma$ is $\gamma\cup\Sigma_{uc}$.
\end{remark}



% Let language $K \subseteq \mathcal{L}_m(G)$ be a non-prefix-closed language that represents the desired behavior to be achieved under control. $K$ is said to be:
% \begin{enumerate}[(i)]
%     \item controllable (w.r.t. DFA $G$ and uncontrollable event set $\Sigma_{uc}$) if $\overline{K} \Sigma_{uc} \cap \mathcal{L}(G) \subseteq \overline{K}$;
%     \item observable (w.r.t. DFA $G$, controllable event set $\Sigma_c$ and observable event set $\Sigma_o$) if for all $s \in \overline{K}$ and $\sigma \in \Sigma_c$,
%     such that $s \, \sigma \in \overline{K}$, $P^{-1}(P(s)) \, \sigma \cap \mathcal{L}(G) \subseteq \overline{K}$, where $P^{-1}(s) := \{ t \in \Sigma^* | P(t) = s \}$;
%     \item $\mathcal{L}_m(G)$-closed if $\overline{K} \cap \mathcal{L}_m(G) = K$. 
% \end{enumerate}
% %
% The following theorem shows that these three conditions together provide the necessary and sufficient conditions for the existence of a supervisor $S$ that exactly achieves the given language $K$.
% %
% \begin{theorem}[Controllability and Observability Theorem \cite{lin1988observability}]
%     Consider DES $G$, with controllable event set $\Sigma_c$ and observable event set $\Sigma_o$. Consider also $K \subseteq \mathcal{L}_m(G)$. 
%     There exists a supervisor $S : \mathcal{L}(G) \rightarrow \Gamma$ such that $\mathcal{L}_m(S /G) = K$ and $L(S / G) = \overline{K}$ if and only if
% \begin{enumerate}[(i)]
%     \item $K$ is controllable w.r.t. $\mathcal{L}(G)$ and $\Sigma_{uc}$;
%     \item $K$ is observable w.r.t. $\mathcal{L}(G)$, $\Sigma_o$ and $\Sigma_c$;
%     \item $K$ is $\mathcal{L}_m(G)$-closed.
% \end{enumerate}
% \end{theorem}

% If a given specification language cannot be exactly achieved,
% then the synthesis problem asks whether we can find a controllable, observable, and $\mathcal{L}_m(G)$-closed sublanguage of the specification language, and one that is ``as large as possible.''
% We would like to synthesize the non-blocking and maximally permissive, that is, the closed-loop system $S / G$ satisfies:
% \begin{enumerate}[(i)]
%     \item $\mathcal{L}_m(S /G) \subseteq K$ and $\mathcal{L}(S /G) = \mathcal{L}_m(S /G)$;
%     \item For any other a supervisor $S'$ satisfying (i), we have $\mathcal{L}_m(S / G) \not \subset \mathcal{L}_m(S' / G)$.
% \end{enumerate}
%
% Note that (ii) is usually relaxed in the scenario under partial observations, since observability is not preserved under union.


\section{Problem formulation}\label{sec:prob}

In this section, we introduce the notion of cognitive non-blocking based on the preliminaries and formulate the synthesis problems for DESs under this notion.
%
% TODO examples? 一个还是两个?
Then, we provide examples to illustrate the challenges of this problem.

\subsection{Cognitive non-blocking synthesis problem under partial observations}

% In the previous section, Definition~\ref{def:nb} introduces the non-blocking property in DESs.
% Under this property, every state that is \emph{accessible} from the initial state $q_0$ is also \emph{co-accessible} to the marked state $Q_m$.
%
The non-blocking supervisory control problem aims to synthesize a supervisor that selects admissible control decisions based solely on partial observations, such that every reachable sequence in the closed-loop system can be extended to reach the marked states. 
In this setting, the supervisor can restrict the behavior of the system by dynamically enabling or disabling the controllable events, thereby rejecting the reachability to the blocking states.

Formally, the non-blocking supervisory control problem under partial observation asks to find a supervisor $S: \mathcal{L}(G) \to \Gamma$ such that
\begin{flalign}\label{eq:nb_sc}
    \forall s \in \mathcal{L}(S / G), \exists t \in \Sigma^* : st \in \mathcal{L}_m(S / G),
\end{flalign}
or, equivalently, in the closed-loop system $S/G$, all states accessible from the initial state are co-accessible to the marked states.
However, non-blockingness only captures the existence of continuations to marked states in the underlying DFA $G$.
In this sense, when the system is only partially observable, and the supervisor does not know the exact current state at each step, reaching a marked state in the actual system may still be insufficient for certifying, based solely on observations, that a marked state has indeed been reached.

Therefore, we introduce a stronger notion, called \emph{cognitive non-blocking}.
Before giving the problem formulation, we first discuss the physical interpretation behind this definition.
Cognitive non-blocking emphasizes the ability to \emph{certify} task completion based on observations.
%
If a marked state is reachable from an unmarked state only through unobservable events, then the system may enter the marked set silently, making it impossible for an external observer to determine whether the task has indeed been completed.
Therefore, interpreting reachability of marked states as task completion is meaningful only when such completion can be certified from observations.
% 
In many physical scenarios, such as a robot arriving at a goal region or a machine finishing an operation, task completion is expected to be accompanied by some observable evidence.
%
Motivated by this observation, we formally introduce the notion of cognitive non-blocking and then formulate the corresponding synthesis problem as follows.
%
\begin{definition}%[Cognitive Non-blocking]
    \label{def:cnb}
    Given DFA $G = (Q, \Sigma, \delta, q_0, Q_m)$ and a set of observable events $\Sigma_o \subseteq \Sigma$, we say DFA $G$ is \textit{cognitive non-blocking} with respect to observable event set $\Sigma_o$ if
    \begin{equation}\label{eq:cnb_1}
    \forall s \in \mathcal{L}(S/G), \exists t \in \Sigma^* : \hat{\mathcal{E}}_S(P (st)) \subseteq Q_m.
    \end{equation}
\end{definition}
%
\begin{problem}%[Cognitive non-blocking synthesis problem]
    \label{prob:cnb}
    Given DFA $G = (Q, \Sigma, \delta, q_0, Q_m)$, a set of observable events $\Sigma_o \subseteq \Sigma$, and a set of controllable events $\Sigma_c \subseteq \Sigma$, synthesize the partial-observation supervisor such that
    \begin{enumerate}[(i)]
        \item The system is \textit{cognitive non-blocking} with respect to  $\Sigma_o$;
        \item The supervisor is maximally permissive in the sense that for any other supervisor $S'$, we have
        $$
            \mathcal{L}(S/G) \not \subset \mathcal{L}(S'/G).
        $$
    \end{enumerate}
\end{problem}
%
Intuitively, cognitive non-blockingness requires that the supervisor can eventually certify, based solely on observations, that the system is in a marked state.
That is, for every current behavior $\alpha$, there exists a continuation $\omega$ such that the resulting immediate state estimate $\hat{\mathcal{E}}_S(\alpha\omega)$ consists entirely of marked states.
%
% Problem~\ref{prob:cnb} is to synthesize a maximally permissive supervisor that enforces cognitive non-blockingness on the closed-loop system.


\subsection{Challenge of the synthesis problem}

\begin{figure}[t]
    \centering
    \includegraphics[width=\linewidth]{Figs/obs_nb_challenge.pdf}
    \caption{The DFAs $G_1$ and $G_2$ are shown in subfigures~(a) and~(b), respectively, with observable event set $\Sigma_o = \{ o_1, o_2, o_3, o_4\}$ and marked-state set $Q_m = \{ q_5 \}$. 
    The shaded regions represent identical observation-induced state estimates, and the corresponding estimate transitions are the same in both systems, showing that $G_1$ and $G_2$ induce the same observer-level structure. 
    $G_1$ is blocking because $q_1$ and $q_3$ are not co-accessible to $q_5$, whereas $G_2$ is non-blocking owing to its cross-coupled transitions.}
    \label{fig:obs_nb_challenge}
\end{figure}

A central challenge is that non-blockingness is not determined by the observer-level structure alone.
An observer captures how state estimates evolve under observable events, but it does not retain the internal transition connectivity among the concrete states in each estimate.
Since non-blockingness requires every reachable state to admit a continuation to a marked state, an observer state may appear co-accessible even though some of its constituent states are blocking.
The following example illustrates this mismatch between the observer-based structure and state-level non-blocking co-accessibility.



\begin{example}\label{ex:obs_nb_challenge_cnb}
Consider the two DFAs $G_1$ and $G_2$ in Fig.~\ref{fig:obs_nb_challenge}, with observable event set
$
    \Sigma_o = \{ o_1, o_2, o_3, o_4 \}
$
and marked-state set
$
    Q_m=\{ q_5 \}.
$
One can easily check that both systems induce the same observation-based estimates $\{ q_1, q_2 \}$ and $\{ q_3, q_4 \}$.
For simplicity, we omit the co-accessibility of states $q_0$ and the intermediate states between $q_0$ and $q_1$, or $q_0$ and $q_2$, and only discuss the non-blocking property of states $q_0$ to $q_5$ in place of the overall system.
%
For DFA $G_1$ of Fig.~\ref{fig:obs_nb_challenge} (a), since states $q_1$ and $q_3$ are not co-accessible to the marked state $q_5$, the system violates the non-blocking property.
In contrast, the cross-coupled transitions in DFA $G_2$ of Fig.~\ref{fig:obs_nb_challenge} (b) ensure that every state can eventually reach $q_5$ via the transition sequence $q_3 \xrightarrow{o_3} q_1 \xrightarrow{o_2} q_4 \xrightarrow{o_3} q_2 \xrightarrow{o_4} q_5$, and hence $G_2$ is non-blocking.
\end{example}

Example~\ref{ex:obs_nb_challenge_cnb} shows that identical observer-level structures may correspond to different non-blocking properties.
We next use the same pair of DFAs to explain why the observer-level structure alone is also insufficient for cognitive non-blocking synthesis under partial observation.

\begin{example}\label{ex:obs_based_structure_insufficient}
Consider again the two DFAs $G_1$ and $G_2$ in Fig.~\ref{fig:obs_nb_challenge}.
Because they induce the same observer-level structure, any synthesis method whose input is restricted to this structure receives identical information for both systems.
%
Restricting attention to states $q_1$ through $q_5$, the observer consists of the states $\{q_1, q_2\}$, $\{q_3, q_4\}$, and $\{q_5\}$ with transitions
$
\{q_1,q_2\} \xrightarrow{o_2} \{q_3,q_4\},
\{q_3,q_4\} \xrightarrow{o_3} \{q_1,q_2\},
\{q_1,q_2\} \xrightarrow{o_4} \{q_5\}.
$
At the observer level, $\{q_1, q_2\}$ appears co-accessible because of the $o_4$-transition enabled by $q_2$.
However, non-blocking synthesis requires that such apparent progress at the estimate level be realizable from every underlying state consistent with the estimate under observation-consistent control decisions.

Classical methods such as backward reachability search~\cite{thuijsman2020transformational,thuijsman2024reducing} or the bipartite all-inclusive controller~\cite{yin2016synthesis} operate solely on partial-observation state estimates or on their aggregated reachability graphs.
%
Consequently, they cannot decide whether individual states such as $q_1$ or $q_3$ should be pruned without additional DFA-state-level transition information.

\end{example}

Examples~\ref{ex:obs_nb_challenge_cnb} and~\ref{ex:obs_based_structure_insufficient} reveal that the observer fails to capture the coupling between observer-level estimation and concrete state-level reachability, which is essential for non-blockingness.
%
Moreover, in cognitive non-blocking synthesis, after every observation-consistent decision history, each state contained in the current observer-level estimate must admit a continuation that is consistent with the induced observer-level evolution and eventually reaches a set of states that can be certified as a subset of the marked states $Q_m$.

This requirement cannot be satisfied using the observer alone, since it discards the internal transition structure within each estimate.
%
Accordingly, Problem~\ref{prob:cnb} requires a structure that simultaneously tracks both the evolution of observation-induced estimates and the underlying DFA-state-level transitions that realize such evolution.
%
Such a structure must explicitly maintain the interaction between these two levels, so that observer-level evolution and state-level reachability can be evaluated in a unified manner.



\section{Bipartite cognitive controller}\label{sec:bcc}

In this section, we attempt to solve the cognitive non-blocking problem using information-state-based structures.
%
We first give the notion of information state and the corresponding update rules to represent both unobservable reach under control patterns and observable reach with respect to specific observations.
%
We then put forward the bipartite cognitive controller (BCC) that captures the microscopic evolution of the information states.


\subsection{Information state and bipartite cognitive controller}

Since we are dealing with partially observed systems, we define the full set of information states as $X = 2^Q$, and use $x, \hat{x} \in X$ to denote the information states.
%
The \emph{unobservable reach} of the information state $x \in X$ under
the control pattern $\gamma \subseteq \Sigma$ is given by
\begin{flalign}
    UR_\gamma(x) := \{ &q' \in Q \mid \exists q \in x, \exists \omega \in ((\Sigma_{uc} \cup \gamma) \cap \Sigma_{uo})^* \cup \notag \\
    & \{ \varepsilon \} : q' = \delta(q, \omega) \}.
    \label{eq:ur}
\end{flalign}
The \emph{observable reach} of information states $x \in X$ under observable event $\sigma_o \in \Sigma_o$ is given by
\begin{equation}
    NX_{\sigma_o}(x):=\{q'\in Q\mid \exists q\in x : \delta(q,\sigma_o)=q'\}.
    \label{eq:nx}
\end{equation}
%

Based on the update rules of information states, we need a structure that can not only track the evolution of the controlled system based on the partial observations, but also track the microscopic structure within unobservable reaches and observable reaches.
%
The bipartite cognitive controller (BCC) structure fulfills this role by expanding information states and recording admissible control sequences.
Now, we give the formal definition of BCC as follows.
%

\begin{definition}\label{def:bcc}
Given a DFA $G$, its corresponding BCC is a tuple
$
\mathfrak{S} = (\hat{\mathbb{I}}, \mathbb{I}, \hat{\imath}_0, \Sigma, \Gamma, \delta_{DO}, \delta_{OO}, \delta_{OD})
Y_m)
$
where
\begin{enumerate}[(i)]
    \item 
    $\hat{\mathbb{I}} \subseteq Q \times 2^Q$ is the set of  \emph{decision states}, 
    where $\hat{\imath}_0 = ( q_0,\{ q_0 \}) \in \hat{\mathbb{I}}$ is the initial state; 
    \item 
    $\mathbb{I}  \subseteq Q \times 2^Q\times \Gamma$ is the set of  \emph{observation states};
    \item 
    $\Sigma$ is the set of events in $G$;
    \item
    $\Gamma \subseteq 2^{\Sigma}$ is the set of control patterns; 
    \item 
    the transition function has three different types: 
    
    $\delta_{DO} : \hat{\mathbb{I}} \times \Gamma \rightarrow \mathbb{I}$ is the partial transition function from decision states to observation states, such that, for any
    $\hat{\imath} = (\hat{q}, \hat{x}) \in \hat{\mathbb{I}}$, 
    $\imath = ({q}, {x},\gamma) \in {\mathbb{I}}$,
    we have $\delta_{DO}(\hat{\imath},\gamma)=\imath$ only if 
    \begin{equation}
        [q=\hat{q}] \wedge [{x} = UR_{\gamma}(\hat{x})];
    \end{equation}
    
    $\delta_{OO} : \mathbb{I} \times \Sigma_{uo} \rightarrow \mathbb{I}$ is the partial transition between observation states, such that, for any
    $\imath = ({q}, {x},\gamma),\imath' = ({q}', {x}',\gamma')  \in {\mathbb{I}},\sigma\in\Sigma_{uo}\cap\gamma$,
    we have $\delta_{OO}({\imath},\sigma)=\imath'$ only if 
    \begin{equation}
        [q'=\delta(q,\sigma)] \wedge [x'=x]  \wedge [\gamma'=\gamma];
    \end{equation}

    \item $\delta_{ZY} : Z \times \Sigma_o \rightarrow Y$ is the partial transition function from observation states to decision states, such that for $y = (x_y, q_y) \in Y$, $z = (x_z, q_z, \gamma_z) \in z$, we have $x_y = NX_{\sigma_o}(x_z)$ for $\sigma_o \in (\Gamma \cup \Sigma_{uc}) \cap \Sigma_o$;

    % TODO
    % \hat{\mathbb{I}}_m 要在后续被定义
    
\end{enumerate}
\end{definition}
%
For each $x\in 2^Q$, we define 
\[
\hat{\mathbb{I}}(x)=\{ (q,x)\in Q\times 2^Q\mid q\in x \}.
\]
For each $x\in 2^Q,\gamma\in \Gamma$, we define 
\[
{\mathbb{I}}(x,\gamma)=\{ (q,x,\gamma)\in Q\times 2^Q\mid q\in x \}.
\]
%
% TODO, 这个用法争取在后续不再使用
%       转为分开定义a run rho ends at control / decision states
%   In BCC $\mathfrak{S}$, we use $P = Y \cup Z$ to denote the set of states.
%
% TODO
% 这个也不再使用
% For a decision state $y = (x_y, q_y)$ and a observation state $z = (x_z, q_z, \gamma_z)$, we further denote
% $X(y) = q_y$, $X(z) = q_z$, $Q(y) = q_y$, $Q(z) = q_z$, and $\Gamma(z) = \gamma_z$.

The run (state event sequence) in BCC $\mathfrak{S}$ is of the following form
\begin{equation}\label{eq:run_bcc}
\begin{aligned}
\rho
={}&
\begin{array}[t]{@{}l@{}}
\hat{\imath}_0
\xrightarrow{\gamma_0} {\imath}_{1,1}
\xrightarrow{\sigma_{u}^{(1,1)}} {\imath}_{1,1}
\xrightarrow{\sigma_{u}^{(1,2)}} {\imath}_{1,2}
\cdots
\xrightarrow{\sigma_{u}^{(1,m_1)}} {\imath}_{1,m_1}
\\[0.6em]
\xrightarrow{\sigma_{o}^{(1)}} \hat{\imath}_1
\xrightarrow{\gamma_1} {\imath}_{2,1}
\xrightarrow{\sigma_{u}^{(2,1)}} {\imath}_{2,1}
\xrightarrow{\sigma_{u}^{(2,2)}} {\imath}_{2,2}
\cdots
\xrightarrow{\sigma_{u}^{(2,m_2)}} {\imath}_{2,m_2}
\\[0.4em]
\multicolumn{1}{c}{\vdots}
\\[0.4em]
\xrightarrow{\sigma_{o}^{(k)}} \hat{\imath}_k
\xrightarrow{\gamma_k} {\imath}_{k,1}
\xrightarrow{\sigma_{u}^{(k,1)}} {\imath}_{k,1}
\xrightarrow{\sigma_{u}^{(k,2)}} {\imath}_{k,2}
\cdots
\xrightarrow{\sigma_{u}^{(k,m_k)}} {\imath}_{k,m_k},
\end{array}
\end{aligned}
\end{equation}
such that the sequences
\begin{flalign}\label{eq:seq_bcc:1}
\rho\restriction_{\imath} &= 
\hat{\imath}_0 {\imath}_{1,1} {\imath}_{1,2} \cdots {\imath}_{1,m_1} \hat{\imath}_1 {\imath}_{1,2}{\imath}_{1,2} \cdots  \notag \\
%
\rho\restriction_{\delta} &=  \gamma_0 \sigma_{u}^{(1,1)} \sigma_{u}^{(1,2)} \cdots {\imath}_{1,m_1} \gamma_1 \sigma_{u}^{(2,1)} \cdots \notag \\
%
\rho\restriction_{X, \Gamma} &= \hat{x}_0 (x_1, \gamma_1) \hat{x}_1 (x_2, \gamma_2) \cdots
\end{flalign}
are for the state sequence, control-event sequence, and macrostate sequences,
and
$\hat{x}_0$, $(x_1, \gamma_1)$, $\hat{x}_1 $, and $(x_2, \gamma_2)$ denote the macrostates of $\hat{\imath}_0 = (q_0, x_0)$, $\imath_{1,j} = (q_{1,j}, x_1, \gamma_1)$, $\hat{\imath}_1 = (\hat{q}_1, \hat{x}_1)$, and $\imath_{2,j} = (q_{2,j}, x_2, \gamma_2)$, respectively.
%
%
% 为什么是o? 因为这是能看见的
% TODO
% follow this idea such that the notations can be simple and instant reasoning
% \[
% \rho\restriction_{o}= 
% \hat{\imath}_0{\imath}_{1,1}\hat{\imath}_1{\imath}_{2,1}
% \dots 
% \hat{\imath}_k{\imath}_{k,1}..?
% \]
%
% TODO sequence of macrostates
% macrostate 不做符号定义，直接x和x_\gamma，但是用这个名字
%
%
% We denote the set of macrostates as $W = W^Y \dot{\cup} W^Z$, where
% $$
% \begin{aligned}
%     W^Y &= \{  x_y \in X \mid (x_y, q_y) \in Y \} \\    
%     W^Z &= \{ (x_z, \gamma_z) \in X \times \Gamma \mid (x_z, q_z, \gamma_z) \in Z \}.   
% \end{aligned}
% $$
% $\hat{\mathbb{I}}\restriction_{X}$
%
% $\mathbb{I}\!\!\restriction_{X\times\Gamma}$
%
% for simplicity, we omit the subscript....
%
%
Intuitively, the BCC structure is introduced to describe the ``game'' between the supervisor and the system under partial observation. 
%
It represents not only the control-observation evolution at the level of information states, but also the microscopic transitions within each information state.
The decision states represent sets of observation-equivalent states and serve as decision-making instances for the supervisor, whereas the observation states represent the resulting information states together with the selected decision $\Gamma(z)$, from which the system resolves the next observable event.
%
The transition from $\hat{\imath} \xrightarrow{\gamma} \imath$ captures the unobservable reach under decision $\gamma$,
The transition $\imath \xrightarrow{\sigma_{uo}} \imath'$ captures the transitions within an unobservable reach,
and the transition $\imath \xrightarrow{\sigma_o} \hat{\imath}'$ captures the observable reach after an event $\sigma_o \in \Sigma_o$.
%
Moreover, to specify the control and evolution of the original DFA, we extract ... from \eqref{eq:run_bcc} and \eqref{eq:seq_bcc:1} as follows,
\begin{flalign}    
    \rho\restriction_{\sigma} &= \sigma_{u}^{(1,1)} \cdots \sigma_{u}^{(1,m_1)} \sigma_{u}^{(2,1)} \cdots \notag \\
    \rho\restriction_{\gamma} &= \gamma_0 \gamma_1 \cdots \gamma_i.
\end{flalign}
%



\begin{algorithm}[t]
\caption{Construction of BCC}
\label{alg:bcc}
\KwIn{DFA $G=(Q,\Sigma,\delta,q_0,Q_m)$ with event partitions
      $\Sigma=\Sigma_c\dot{\cup}\Sigma_{uc}=\Sigma_o \dot{\cup} \Sigma_{uo}$\;}
\KwOut{BCC $\mathfrak{S} = (\hat{\mathbb{I}}, \mathbb{I}, \hat{\imath}_0, \Sigma, \Gamma, \delta_{DO}, \delta_{OO}, \delta_{OD})$\;}

Initialize $\Gamma = 2^{\Sigma_c}$, $y_0 = (q_0, \{q_0\})$\;

\ForEach{unprocessed decision macrostate $\hat{x}$}{                   % TODO
    Include all $(\hat{q}, \hat{x})$, $\forall \hat{q} \in \hat{x}$ into $\hat{\mathcal{I}}$\;

    % \tcp{Enumerate all control patterns for all-inclusive construction}
    \ForEach{control pattern $\gamma \in \Gamma$}{
        Generate observation macrostate $(x, \gamma)$, for $x \gets UR_{\gamma}(\hat{x})$\;
        Include all observation states $(q, x, \gamma)$, $\forall q \in x$\;
        Add decision-observation transitions $(\hat{q}, \hat{x}) \xrightarrow[]{\gamma} (q, x, \gamma)$, $\forall q_y = q_z$\;  % 就算$q_y$有自循环也需要这么做
        Add observation-observation transitions $(q, x, \gamma) \xrightarrow[]{\sigma_{uo}} (q', x, \gamma)$, $\forall q \xrightarrow[]{\sigma_{uo}} q'$\;

        \ForEach{observable event $\sigma_o \in (\gamma \cup \Sigma_{uc})\cap\Sigma_o$}{    % TODO 这个说法改一下, 我觉得前面的remark得留下来说明这个事, 不然一直很别扭
            Compute $\hat{x}' = NX_{\sigma_o}(x)$\;
            Include $(\hat{q}', \hat{x}')$ $\forall \hat{q}' \in \hat{x}'$ and add observation-decision transitions $(q, x, \gamma) \xrightarrow[]{\sigma_o} (\hat{q}', \hat{x}')$\;

            Enqueue $w_y' = x_y'$ to $W^Y$ if it has not been processed\;
        }
    }
}

% Set $\mathbb{I}_m = \{ (q, x) \in \mathbb{I} \mid x \subseteq Q_m \}$ and $\Sigma_{\mathfrak{S}} \gets \Sigma \cup \Gamma$\;

\Return $\mathfrak{S}$\;
\end{algorithm}
%
We give Algorithm~\ref{alg:bcc} for the construction of a BCC from the original DFA $G$.
Algorithm~\ref{alg:bcc} follows a breadth-first search (BFS) structure.
It iteratively computes all successor observation macrostates from each decision macrostate under all admissible control patterns, and then derives successor decision macrostates from the observable projections of the newly generated observation macrostates.
Meanwhile, whenever a new macrostate is generated, the algorithm also constructs the associated decision and observation states, together with the transition relations induced by the underlying DFA transitions.
%
In Lines 2--7, the algorithm explores decision macrostates in a BFS manner.
For each decision macrostate $\hat{x}$, it adds all corresponding decision states $(\hat{q}, \hat{x})$ with $\hat{q} \in \hat{x}$ into the BCC, and enumerates all control patterns to compute successor observation macrostates via the operator $UR(\cdot)$.
% The corresponding decision-observation micro-transitions, which correspond to the selected control decision, are then added to the BCC in Line 7.
%
In Lines 8, the algorithm expands each newly generated observation macrostate by adding all observation states $(q, x, \gamma)$ and recording the internal unobservable transitions within the same observation macrostate.
%
In Lines 9--12, the algorithm computes successor decision macrostates using the operator $NX(\cdot)$ for observable events, and adds the corresponding decision states and observation-decision transitions to the BCC.
%
We use Example~\ref{ex:bcc} to illustrate the construction of BCC.
%
\begin{example}\label{ex:bcc}
    Let us consider again DFA $G$ in Fig.~\ref{fig:dfa_grid_env} (b) with observable event set $\Sigma_o = \{o_1, o_2, o_3\}$ and controllable event set $\Sigma_c = \{ c_1, c_2, o_3 \}$. 
    We construct the corresponding BCC using Algorithm~\ref{alg:bcc}, with the resulting structure shown in Fig.~\ref{fig:bcc_overall}, where orange rectangles and blue rounded rectangles denote decision macrostates and observation macrostates, respectively. 
    For clarity, we simplify the notation $(q, \hat{x})$ or $(q, x, \gamma)$ to $q$ and use different colors for the transition relations to highlight distinct evolutions, thereby illustrating the structural properties of each state.

    We also demonstrate the evolution of the BCC as follows.
    %
    For the initial macrostate $w_{y,0} = \{ q_0 \}$, since $u_1, u_2$ are uncontrollable and unobservable, the only available control pattern is $\gamma_0 = \{ o_3 \}$.
    Next, we have $x_{1} = \{ q_0, q_1, q_2 \} = UR_{\gamma_0}(\hat{\imath}_0)$, and the corresponding observation state set is  
    $
    (q_0, x_{1}, \gamma_1),
    $
    $
    (q_1, x_{1}, \gamma_1),
    $
    and
    $(q_2, x_{1}, \gamma_1)$,
    where states $q_0$, $q_1$, $q_2$ are accessible from state $q_0$ through the unobservable transitions $\varepsilon$, $u_1$, and $u_2$, respectively.
    %
    For observation macrostate $w_{z,1}$, the observable reach based on observation $o_3$ and $o_2$ is
    $x_{y,1} = NX_{o_3}(x_{z,1}) = \{ q_3, q_4 \}$ and $x_{y,2} = NX_{o_2}(x_{z,1}) = \{ q_5, q_6 \},$ 
    and the corresponding decision state set is 
    $$
    Y_{x_{y,1}, \gamma_1} = \{(x_{y,1}, q_3), (x_{y,1}, q_4) \}.
    $$
    %
    Similarly, based on the possible unobservable transition $q_3 \xrightarrow{c_1} q_7 \xrightarrow{c_2} q_9$ and $q_4 \xrightarrow{c_1} q_8 \xrightarrow{c_2} q_{10}$, we can enable no event, namely $\emptyset$, or enable $\{ c_1 \}$, $\{ c_2 \}$, or $\{ c_1,c_2 \}$ at this step.
    Then, we have
    $x_{z,4} = UR_{\gamma_2}(x_{y,1}) = \{ q_3, q_4, q_7, q_8 \}$ and
    $x_{z,7} = UR_{\gamma_3}(x_{y,1}) = \{ q_3, q_4, q_7, q_8, q_9, q_{10} \}$,
    where $\gamma_2 = \{ c_1 \}$ and $\gamma_3 = \{ c_1, c_2 \}$.
    For observation macrostates $w_{z,7} = (x_{z,7}, \gamma_3)$, we have
    $x_{y,5} = NX_{o_2}(x_{z,7}) = \{ q_{13}, q_{14} \}$ and
    $x_{y,3} = NX_{o_1}(x_{z,7}) = \{ q_{17}, q_{18} \}$.
\end{example}

\begin{figure*}[t!]
    \centering
    \includegraphics[width=0.615\linewidth]{Figs/bcc_overall_2.pdf}
    \caption{Bipartite cognitive controller constructed from the DFA $G$ in Fig.~\ref{fig:dfa_grid_env} (b). We simplify the notations $(x_y, q_y)$ and $(x_z, q_z, \gamma_z)$ to $q_y$ and $q_z$ to better illustrate the structural properties of each state. 
    Orange rectangles represent decision macrostates, while blue rounded rectangles represent observation macrostates. 
    Transitions in different colors are used to distinguish different evolutions. Moreover, the information states marked by red crosses are infeasible observation macrostates, cognitively blocking decision macrostates, or inconsistent decision or observation macrostates, which should be pruned to achieve cognitive non-blocking. 
    The source code is available at \url{https://github.com/trigger1996/des_common}.}   % TODO, 最好拆成几个库
    \label{fig:bcc_overall}
\end{figure*}


\subsection{Properties of bipartite cognitive controller}


The BCC represents the system evolution at two levels: the macrostate level and the DFA-state level. 
At the macrostate level, it captures the evolution induced by partial observations and control decisions. 
At the DFA-state level, it records the detailed transitions among individual DFA states, including internal transitions within unobservable reaches and observable transitions between decision or observation states.
%
For BCC $\mathfrak{S}$, define
\begin{flalign}\label{eq:pre_post}
    Pre(p)  = \{ &p' \in P \mid \exists \eta \in \Sigma_{\mathfrak{S}}, p \in \delta_{\mathfrak{S}}(p', \eta) \} \nonumber \\
    %
    Post(p) = \{ &p' \in P \mid \exists \eta \in \Sigma_{\mathfrak{S}}, p' \in \delta_{\mathfrak{S}}(p, \eta) \}
\end{flalign}
as the sets of predecessors and successors of $p \in P = Y \ \cup Z$, respectively. 
%
For a set of states $P_{s} \subseteq P$, we use
\begin{flalign}\label{eq:pre_post:2}
    Pre(P_{s})  := &\bigcup_{p \in P_{s}} Pre(p) \nonumber \\
    Post(P_{s}) := &\bigcup_{p \in P_{s}} Post(p)
\end{flalign}
to denote the sets of predecessors and successors, respectively. 
%
For $p, p' \in P$ and $P_s \subseteq P$, we write $p' \leadsto_{P} p$ if there exists a nonempty path from $p'$ to $p$,
$$
\rho_{\mathfrak{S}} = p' p_1 p_2 \cdots p_n p,
$$
such that $p_1, \ldots, p_n \in P_s$.
A state $p \in P$ is called a \emph{deadlock} if 
\begin{equation}\label{eq:bcc_dl}
    Post(p) = \emptyset.
\end{equation}
A nonempty set $P_l \subseteq P$ is called a \emph{livelock} if $P_l$ is a nontrivial strongly connected component, and $Post(P_l) \subseteq P_l$, i.e.,
\begin{flalign}\label{eq:bcc_ll}
    &P_l \neq \emptyset \wedge
    \forall p, (\bm{q}',q') \in P_l : 
    (p \leadsto_{P_l} p') \ \wedge \nonumber \\
    &  (p' \leadsto_{P_l} p) \ \wedge \ Post(P_l) \subseteq P_l.   
\end{flalign}










Based on the above notations, we analyze the properties of the BCC, illustrating the consistency of its evolution with the original DFA, the completeness of the control representation, and its construction complexity in terms of both time and space.
%
% We first clarify a structural restriction on decision macrostates. 
Recall that decision macrostates represent observation-equivalent states under supervisory control, and thus transitions between decision states must preserve the observation structure. In particular, for any $w_y = x_y \in W^Y$ and any $y, y' \in Y_{x_y}$, we have
$
(y, \eta, y') \notin \delta_{\mathfrak{S}}, \forall \eta \in \Sigma_{\mathfrak{S}}.
$
%
Building on this restriction, we describe the internal structure of observation macrostates $w_z = (x_z, \gamma_z)$ and their corresponding observation state $z \in Z_{x_z}$.
A observation state is called an \emph{entry observation state} if all its predecessors $Pre(z)$ are decision states, and if $Post(z) \cap Y \neq \emptyset$, it is called an \emph{exit observation state}. 
Note that a observation state can be both an entry and an exit state.
%
For an entry observation state $z = (x_z, q, \gamma)$ and any predecessor decision state $y = (x_y, q)$, since $\varepsilon$ is admissible to be the transition from $y$ to $z$ at the DFA state level, we have $q \in x_y \cap x_z$.
% This shows that observation states inherit their information state from decision states under unobservable transitions.
%
Moreover, transitions within a observation macrostate are induced by an unobservable reach $UR(\cdot)$ starting from its predecessor decision macrostates.
%
% Therefore, reachability within a observation state is consistent with this inheritance structure. 
Specifically, for any entry observation state $z = (x_z, q, \gamma)$ and its predecessor decision state $y$, every $q' \in x_z$ is reachable via a string 
$
s \in ((\Gamma(z)\cup\Sigma_{uc})\cap\Sigma_{uo})^*
$ 
such that
$$
(x_z, q_y, \gamma) = \delta_{\mathfrak{S}}((x_y, q_y), \gamma),
\quad
(x_z, q_y', \gamma) = \delta_{\mathfrak{S}}((x_z, q_y, \gamma), s).
$$
%
We give the following proposition to illustrate the consistency between the DFA $G$ and the BCC $\mathfrak{S}$.
%
\begin{proposition}\label{prop:bcc_dfa_consistency}
Given a run 
$$
\begin{aligned}
\rho_{\mathfrak{S}} = &y_0 \, z_{1, 1} \, \cdots \, z_{1, m_1} \, y_1 \, y_i \, z_{i+1, 1} \, \cdots \, z_{i+1, m_{i+1}} \ y_{i+1} \cdots
\end{aligned}
$$
in the BCC $\mathfrak{S}$, there exists a corresponding sequence
% 这个s是对的, 因为这是 DFA 序列
$
s = \xi_0 \, \sigma_{o,0} \, \xi_1 \, \sigma_{o,1} \, \cdots \, \xi_i \, \sigma_{o,i} \cdots
$
in the closed-loop system $S/G$ such that:
\begin{enumerate}[(i)]
    \item Each $\xi_i \in ((\Gamma(z_i) \cup \Sigma_{uc}) \cap \Sigma_{uo})^*$ is an unobservable subsequence corresponding to the unobservable reach $z_{i,1} \cdots z_{i,m_{i}}$,
    \item For each subsequence $z_{i,1} \cdots z_{i,m_i}$, there exist events $\sigma_{i,j} \in \Sigma_{uo}$ such that
    $
    (q_{i,j}, \sigma_{i,j}, q_{i,j+1}) \in \delta,
    $
    \item For each transition $z_{i,m_i} \xrightarrow[]{\sigma_{o,i}} y_{i+1}$, we have
    $
    (q_{i,m_i}, \sigma_{o,i}, q_{i+1}) \in \delta.
    $
\end{enumerate}
\end{proposition}
%
\begin{proof}
The result follows directly from the construction of $\delta_{\mathfrak{S}}$.
%
First, every observation-observation transition satisfying $z_{i,j+1} = \delta_{\mathfrak{S}}(z_{i,j}, \sigma)$, $i,j \in \mathbb{N}$ is defined only if $\sigma \in (\Gamma(z) \cup \Sigma_{uc}) \cap \Sigma_{uo}$ and $(q_{i,j}, \sigma, q_{i,j+1}) \in \delta$. 
Thus, each subsequence $z_{i,1} \cdots z_{i,m_i}$ corresponds to a sequence $\xi_i \in ((\Gamma(z_i) \cup \Sigma_{uc}) \cap \Sigma_{uo})^*$ satisfying (ii).
%
Second, every observation-decision transition $z_{i,j} \xrightarrow{\sigma_o} y_{i+1}$ is defined only if $\sigma_o \in \Sigma_o$ and $(q, \sigma_o, q') \in \delta$ of DFA $G$, which establishes (iii).
%
Finally, every decision-observation transition $y_i \xrightarrow{\gamma_{i+1}} z_{i+1,1}$ is defined by $\delta_\mathfrak{S}$ such that $\gamma_{i+1} = \Gamma(z_{i+1})$.
This determines the set of admissible events $(\Gamma(z_{i+1}) \cup \Sigma_{uc})$ that govern the subsequent unobservable evolution within $z_{i+1}$, and thus constrains the corresponding subsequence $\xi_{i+1}$.
Concatenating all such unobservable subsequences $\xi_i$ and observable events $\sigma_{o,i}$ yields a valid execution $s$ in the closed-loop system $S / G$, which proves (i).
\end{proof}
%
As a direct consequence of Proposition~\ref{prop:bcc_dfa_consistency}, 
for any $\alpha \in \mathcal{L}(\mathfrak{S})$, if
$
\delta_{\mathfrak{S}}(y_0,\alpha)=(\bm{q},q),
$
then
$
    \delta(q_0,P_\Sigma(\alpha)) = q.
$
In particular, we have
$$
    \mathcal{L}_{\Sigma}(\mathfrak{S})
    =
    P_\Sigma(\mathcal{L}(\mathfrak{S}))
    \subseteq
    \mathcal{L}(G).
$$
%
This means that every event sequence generated by the BCC and projected onto the DFA alphabet is feasible in $G$. Hence, the BCC does not introduce any spurious DFA behavior; it only refines feasible executions of $G$ by explicitly recording the interaction among observation, control patterns, and unobservable evolution.

We next summarize the computational cost of constructing the BCC. 
Since Algorithm~\ref{alg:bcc} explores decision macrostates in a breadth-first-search manner, the number of possible decision macrostates is bounded by $2^{|Q|}$. 
For each generated decision macrostate, the algorithm enumerates control patterns $\gamma \in 2^{\Sigma_c}$. 
For each pair of a decision macrostate and a control pattern, the computation of the corresponding unobservable reach, internal observation-observation transitions, and observable successors can be implemented by scanning the DFA transition table, whose cost is bounded by $O(|Q|\cdot |\Sigma|)$. Therefore, the worst-case time complexity of Algorithm~\ref{alg:bcc} is
$$
O(2^{|Q| + |\Sigma|} \cdot |Q| \cdot |\Sigma|)
$$
since $|\Sigma_c|\leq |\Sigma|$.

The size of the constructed BCC is also finite. 
Each BCC state is either a decision state of the form $(x,q)$ or a observation state of the form $(x,q,\gamma)$. Since a macrostate $x$ is a subset of $Q$ and a control pattern satisfies $\gamma\in 2^{\Sigma_c}$, the total number of BCC states satisfies
\begin{equation}
    \label{eq:prop:bcc:1}
    |Y + Z| \leq |Q|(2^{|Q|}+2^{|Q|+|\Sigma|}).    
\end{equation}
Moreover, the transitions of the BCC consist of decision-observation transitions, internal observation-observation transitions, and observation-decision transitions. Since each observation state has at most one successor under each event of the deterministic DFA $G$, the total number of BCC transitions is bounded by
\begin{equation}
    \label{eq:prop:bcc:2}
    |\delta_{\mathfrak{S}}| \leq |Q|\,2^{|Q|+|\Sigma|}(1+|\Sigma|).
\end{equation}
Thus, Algorithm~\ref{alg:bcc} constructs a finite BCC in finite time, and both its state space and transition relation are bounded exponentially in the number of DFA states and controllable-event choices.

% That is, every DFA sequence projected from the BCC is feasible and belongs to the language $\mathcal{L}(G)$.
% %
% Moreover, we provide the time complexity of Algorithm~\ref{alg:bcc} and the space complexity of BCC as follows.

% \begin{proposition}
% Given a DFA $G$ with $|Q|$ states and $|\Sigma|$ events, the complexity of Algorithm~\ref{alg:bcc} in time will be dominated by
% $$
% O(2^{|Q| +|\Sigma|} \cdot (|Q|\cdot |\Sigma|))
% $$
% \end{proposition}
%
% \begin{proof}
% Algorithm~\ref{alg:bcc} performs a breadth-first search over generated decision macrostates. Since each decision macrostate is a subset of $2^Q$, the number of repetitions of Line 2 is at most
% $
%     |W^Y| \leq 2^{|Q|}.
% $
% %
% In Line 3, the decision states will be added to BCC, and the repetition will be at most  $O(|Q|)$.
% %
% In Line 4, for each generated decision macrostate $w_y$, the algorithm enumerates all control patterns
% $
%     \gamma \in \Gamma=2^{\Sigma_c},
% $
% such that there are at most
% $
%     |\Gamma| = 2^{|\Sigma_c|}
% $
% iterations of the inner loop.
% %
% We next estimate the cost for a fixed pair $(w_y, \gamma)$ in Line 5, which corresponds to the admissible control pattern at state $w_y$.
% Denote the set of admissible observable and unobservable reachable events as
% $
%     \mathcal{U}_\gamma = (\gamma \cup \Sigma_{uc})\cap\Sigma_{uo},
% $
% and
% $
%     \mathcal{O}_\gamma = (\gamma \cup \Sigma_{uc})\cap\Sigma_o .
% $
% First, the algorithm computes
% $
%     x = UR_\gamma(y).
% $
% Using a standard graph-search implementation over the DFA transition table, this requires visiting at most $|Q|$ states and checking at most 
% $| \mathcal{U}_\gamma | \leq |\Sigma|$ events at each state. Thus, the computation complexity of $UR_\gamma(y)$ is
% $
%     O(|Q|\cdot |\Sigma|).
% $
% %
% After $x_z = X(z)$ is obtained, the algorithm adds the corresponding observation state and records the internal unobservable transitions inside the same observation macrostate. 
% This amounts to scanning the subgraph induced by $x$ under events in $\mathcal{U}_\gamma$.
% Since $|x_z|\leq |Q|$ and the number of such internal transitions is at most
% $|Q|\cdot |\Sigma|$, this step costs
% $
%     O(|Q|\cdot |\Sigma|).
% $
% %
% Finally, in Line 9, for each $\sigma_o \in \mathcal{O}_\gamma$, it computes
% $
%     y' = x_y' = NX_{\sigma_o}(x)
% $
% and records the corresponding observation-decision micro-transitions. 
% Since $|\mathcal{O}_\gamma| \leq |\Sigma|$ and each computation scans at most $|Q|$ states, this step also costs
% $
%     O(|Q|\cdot |\Sigma|).
% $
% The total cost of Algorithm~\ref{alg:bcc} is
% $$
% O(|W^Y| \cdot 2^{|\Sigma|} \cdot (|Q| + 3(|Q|\cdot |\Sigma|))),
% $$
% which will be dominated by
% $
% O(2^{|Q| +|\Sigma|} \cdot (|Q|\cdot |\Sigma|)).
% $
% \end{proof}



% \begin{proposition}\label{prop:state_bcc}
% Consider a DFA $G$ with $|Q|$ states and $|\Sigma|$ events. For the corresponding BCC $\mathfrak{S}$ constructed by Algorithm~\ref{alg:bcc}, the number of states and transitions satisfy
% $$
% \begin{aligned}
%     |Y + Z| &\leq |Q|(2^{|Q|}+2^{|Q|+|\Sigma|}) \\
%     |\delta_{\mathfrak{S}}| &\leq |Q|\,2^{|Q|+|\Sigma|}(1+|\Sigma|),
% \end{aligned}
% $$
% respectively.
% \end{proposition}
%
% \begin{proof}
% Each BCC state is of the form $(x, q)$ or $(x, q, \gamma)$, so the maximal number of states satisfies
% $
%     |Y + Z| \leq (|W^Y| \cdot |Q| + |W^Z| \cdot |Q| \cdot |\Gamma|).
% $
% Since each decision state is a subset of $Q$, we have
% $
%     |W^Y| \leq 2^{|Q|}.
% $
% As for the observation macrostates in the form of
% $
%     w_z = (x, \gamma),
% $
% where $x \subseteq 2^Q$ and $\gamma\in 2^{\Sigma_c}$. 
% Therefore, the maximal number of observation macrostates is
% $
%     |W^Z| \leq 2^{|Q|}\cdot 2^{|\Sigma_c|} = 2^{|Q|+|\Sigma_c|},
% $
% and it follows that
% $$
%     |Y + Z| \leq |Q|(2^{|Q|}+2^{|Q|+|\Sigma_c|}).
% $$

% Next, we calculate the maximal number of transitions. 
% The BCC contains three types of micro-transitions: decision-observation transitions, internal observation-observation transitions, and observation-decision transitions.
% %
% For decision-observation transitions, for each decision macrostate $y$, the control pattern satisfies $\gamma \in 2^{\Sigma_c}$.
% Thus, the number of decision-observation transitions $(x_y, q_y) \xrightarrow{\gamma} (x_z, q_z, \gamma)$ is bounded by
% $$
%     |W^Y|\cdot 2^{|\Sigma_c|} \leq |Q|\,2^{|Q|+|\Sigma_c|}.
% $$
% %
% Given a observation macrostate $w_z = (x_z, \gamma)$, each corresponding observation state $(x_z, q_z, \gamma)$ has at most one successor under each event $\sigma \in \Sigma$, since the DFA $G$ is deterministic.
% Therefore, the total number of internal observation-observation transitions and observation-decision transitions over all observation macrostates is bounded by
% $
%     |W^Z| \cdot |Q|\cdot|\Sigma|.
% $
% Since
% $
%     |W^Z| \leq 2^{|Q|+|\Sigma_c|},
% $
% we obtain
% $
%     |W^Z| \cdot |Q|\cdot|\Sigma| \leq |Q| \cdot |\Sigma| \cdot 2^{|Q|+|\Sigma_c|}.
% $
% %
% Combining the above bounds, we obtain
% $$
% \begin{aligned}
%     |\delta_{\mathfrak{S}}|
%     &\leq
%     |Q|\,2^{|Q|+|\Sigma_c|}
%     +
%     |Q|\cdot|\Sigma|\cdot 2^{|Q|+|\Sigma_c|} \\
%     &=
%     |Q|\,2^{|Q|+|\Sigma_c|}(1+|\Sigma|).
% \end{aligned}
% $$
% %
% Since $|\Sigma_c| \leq |\Sigma|$, we further have
% $$
%     |Y + Z| \leq |Q|(2^{|Q|}+2^{|Q|+|\Sigma|}),
% $$
% and
% $$
%     |\delta_{\mathfrak{S}}| \leq |Q|\,2^{|Q|+|\Sigma|}(1+|\Sigma|).
% $$
% \end{proof}




\section{Solution to Cognitive non-blocking synthesis problem}\label{sec:solu}
% 
% 我后面需要什么，我前面需要什么
% 
%
% 1 consistency of BCC 作为下一章节的第一个定理
% 2 what is potentially feasible observation state?
% 3 what is cognitive live decision state
% 4 how to synthesis CNB
% 5 Algorithm 2
% 6 correctness of Algorithm 2 (有必要用单源单汇路径的结论吗？)
% 7 

In this section, we leverage the definition of consistency for BCCs.
By pruning those inconsistent states in BCC, the sublanguage generated by the BCC can be guaranteed to be cognitively non-blocking.
%

\subsection{Synthesis algorithm}\label{subsec:supervisor_induction}

To synthesize cognitive non-blocking sublanguages, we first identify the target BCC states characterized by \eqref{eq:cnb_1}.
%
According to \eqref{eq:cnb_1}, given a cognitively marked decision marcostate, all its decision state are marked. 
Therefore, regardless of which actual DFA state has been reached under the partial observation, it can be certified that the system has reached a marked state.
%
Meanwhile, we define the \emph{consistency} of a BCC $\mathfrak{S}$ as follows.
%
\begin{definition}\label{def:con_bcc}
Given a BCC $\mathfrak{S}$, we say it is \emph{consistent} if
\begin{enumerate}[(i)]
    \item For each observation state $(x_z, q_z, \gamma_z)$, all possible observable transitions from its marcostate $w_z = (x_z, \gamma_z)$ are defined, 
    $
    \forall (x_z, q_z, \gamma_z) \in Z_{x_z, \gamma_z}, \forall \sigma_o \in \Gamma(z)\cap\Sigma_o:\;
    (\delta(q_z,\sigma_o)!) \Rightarrow \delta_{\mathfrak{S}}((x_z, q_z, \gamma_z), \sigma_o)!;    
    $
    \item For each decision macrostate $w_y = x_y \in W^Y$, at least one control decision is defined for its decision state $y = (x_y, q_y) \in Y_{x_y}$, 
    $
    \forall (x_y, q_y) \in Y_{x_y}, w_y = x_y  : \mathcal{H}(w_y) \neq \emptyset.
    $
\end{enumerate}
\end{definition}
%
Note that in Definition~\ref{def:bcc} (iii), only observation states can be admissible successors for decision states in BCC $\mathfrak{S}$.
%
Next, we give the following two definitions for properties of decision macrostates and observation macrostates based on the reachability of their corresponding decision or observation states, which will be further utilized for the notion of cognitive non-blocking.
%
% \begin{definition}\label{def:z_pf}
%     We say a observation macrostate $z \in Z$ is \emph{potentially feasible} if, for each of its states $q^Z \in \bm{\mathcal{Q}}_z$, it does not lead to a deadlock \eqref{eq:bcc_dl} or a livelock \eqref{eq:bcc_ll} consisting solely of the observation states $\bm{\mathcal{Q}}_z$ of $z$. 
% \end{definition}
\begin{definition}\label{def:z_pf}
    We say a observation macrostate $w_z = (x_z, \gamma_z)$ is \emph{potentially feasible} if, for each of its observation states $z \in Z_{x_z}$, there does not exist a deadlock \eqref{eq:bcc_dl} from $z$, nor does there exist a livelock \eqref{eq:bcc_ll} contained entirely in $w_z$.
\end{definition}
%
% 这里有一点很关键，这个decision state必须能到同一个mark，不能(y, q)到一个marked, (y, q')到另一个mark
\begin{definition}\label{def:y_cl}
    We say a decision macrostate $w_y = x_y$ is \emph{cognitively live} with respect to the set of marked decision states $Y_m$ if and only if there exists a cognitively marked decision state $y_m = (x_{y,m}, q_{y,m}) \in Y_m$ and an information-state sequence
    $
    \iota = w_y w_{z,0} w_{y,1} w_{z,1} \cdots w_{z,k} y_m,
    $
    such that all decision states $y = (w_y, q_y)$ which correspond to $w_y$ satisfy
    \begin{flalign}\label{eq:bcc_y_cl}
        &\forall q_y \in x_y, \exists q_{y,m} \in x_{y,m} : \rho_{\mathfrak{S}} \in R_X^{-1}(\iota) : \notag \\
        &y_m = \delta_{\mathfrak{S}}(y, tr_{\mathfrak{S}}(\rho_{\mathfrak{S}})),
    \end{flalign}
    where $w_{z,k}$, $w_{z,k'}$, $k \in \mathbb{N}$, $k' \in \mathbb{N}_{\geq 0}$ are potentially feasible observation macrostates and other cognitively live decision macrostates, respectively.
\end{definition}
%
Intuitively, a decision macrostate is cognitively live if all its decision states can reach the decision states of the other cognitively marked decision macrostate via the runs $R_X^{-1}(\iota)$ corresponding to the sequence of macrostates, $\iota$, of the potentially feasible observation macrostates and other cognitively live decision macrostates.

The decision and observation state-level of Definition~\ref{def:z_pf} and \ref{def:y_cl} are defined upon their corresponding information states.
If a observation macrostate is potentially feasible, then all of its observation states are potentially feasible, and vice versa.
If a decision macrostate is said to be cognitively live, then all its decision states are cognitively live.
%
% 第二个没有 vice versa
%
We give the following example to illustrate in detail the property of potential feasibility and cognitive liveness in a BCC.


%
\begin{example}\label{ex:y_z_prop}
Let us consider again the DFA $G$ in Fig.~\ref{fig:dfa_grid_env} (b) and the corresponding BCC $\mathfrak{S}$ in Fig.~\ref{fig:bcc_overall}.
%
For the branch of $w_{y,0}, w_{z,1}, w_{y,1}, w_{z,7}, w_{y,5}, w_{z,6}, w_{y,3}, w_{z,9}$, the observation macrostate $w_{z,6} = (\{ q_{13}, q_{14} \}, \emptyset)$ is not potentially feasible, since both of its observation state $(\{ q_{13}, q_{14} \}, q_{13}, \emptyset)$ and $(\{ q_{13}, q_{14} \}, q_{14}, \emptyset)$ are deadlock states.
%
Meanwhile, the decision macrostate $w_{y,5} = \{ q_{13}, q_{14} \}$ is not cognitively live, because its only available transition leads to the deadlock observation macrostate $w_{z,6}$.
The decision macrostates $w_{y,1} = \{ q_{3}, q_{4} \}$ and $w_{y,3} = \{ q_{17}, q_{18} \}$ are not cognitively live either.
Since the only available marked decision macrostate is $w_{z,2} = \{ q_5, q_6 \}$, there is no available sequence from the decision state  $(x_{y,1}, q_3)$, $(x_{y,1}, q_4)$, $(x_{y,3}, q_{17})$, or $(x_{y,3}, q_{18})$ to the decision states $(x_{y,2}, q_5)$ and $(x_{y,2}, q_6)$, where $x_{y,1} = \{ q_3, q_4 \}$, $x_{y,2} = \{ q_5, q_6 \}$, and $x_{y,3} = \{ q_{17}, q_{18} \}$. 
Thus, \eqref{eq:bcc_y_cl} is not satisfied.


As for the branch of $w_{y,0}, w_{z,0}, w_{y,2}, w_{z,5}, w_{y,4}, w_{z,12}, w_{z,11}$, the decision macrostate $w_{y,4} = x_{y,4} = \{ q_{15}, q_{16} \}$ is cognitively live.
For the state sequence $\iota = w_{y,4} w_{z,11} w_{y,0} w_{z,0}$, with $w_{z,11}  = (x_{z,11}, \gamma_{z, 11}) = (\{ q_{15}, q_{16}, q_{17}\}, \{c_1, o_3\})$, $w_{y,0} = x_{y,0} = {q_0}$, and $w_{z,0} = (x_{z,0}, \gamma_{z, 0}) = (\{ q_0, q_1, q_2 \}, \emptyset)$, we have the corresponding BCC runs
$$
\begin{aligned}
\rho_{\mathfrak{S}, 1} = 
&(x_{y,4}, q_{15}) \xrightarrow{\{c_1, o_3\}} (x_{z,11}, q_{15}, \gamma_{z,11}) \xrightarrow{c_1} (x_{z,11}, q_{17}, \gamma_{z, 11}) \xrightarrow{o_3} \\
&(x_{y,0}, q_0) \xrightarrow{\emptyset} (x_{z_0}, q_0, \gamma_{z,0}) \xrightarrow{u_1} (x_{z_0}, q_1, \gamma_{z,0}) \xrightarrow{o_2} (x_{y,2}, q_5) \\
% \end{aligned}
% $$
% and
% $$
% \begin{aligned}
\rho_{\mathfrak{S}, 2} = 
&(x_{y,4}, q_{16}) \xrightarrow{\{c_1, o_3\}} (x_{z,11}, q_{16}, \gamma_{z, 11}) \xrightarrow{o_1} (x_{y,0}, q_0) \xrightarrow{\emptyset} \\
&(x_{z,0}, q_0, \gamma_{z,0}) \xrightarrow{u_2} (x_{z,0}, q_2, \gamma_{z, 0}) \xrightarrow{o_2} (x_{y,2}, q_6),
\end{aligned}
$$
both of which end at decision states associated with the same cognitively marked decision state $w_{y,2}$.
%
In this sense, one can also check that $(w_{y,0}, q_0)$ is cognitively live.
Moreover, the observation macrostates $w_{z,11}$ and $w_{z,5} = (x_{z,5}, \gamma_{z,5}) = (\{q_5, q_6, q_{11}, q_{12}\}, \{ c_1 \})$ are potentially feasible, since their exit observation state, $(x_{z,5}, q_{11}, \gamma_{z,5})$, $(x_{z,5}, q_{12}, \gamma_{z,5})$, $(x_{z,11}, q_{16}, \gamma_{z,11})$, and $(x_{z,11}, q_{17}, \gamma_{z,11})$, have cognitively live successors.
However, if the controllable event $o_3$ is not enabled at $w_{y,4}$, then the run $\rho_{\mathfrak{S}, 1}$ is no longer feasible.
Therefore, $w_{y,4}$ is no longer cognitively live.
%  
% That is, by observing the subsequence $(o_3 o_2)^\omega$ or $(o_1 o_2)^\omega$, the observer can certify that the system arrives at a marked state.
%
% 这个问题在于15、16到q_0的trace不一样，一个o_1，一个o_3, 所以不能是run到run，必须是state sequence到run
%
\end{example}


%
Example~\ref{ex:y_z_prop} shows that the macostates must be retained or removed as whole units. 
Since these macrostates represent observation-equivalent DFA states, retaining only part of its states would break the observation-based consistency. 
% Thus, if any associated states violates Definition~\ref{def:con_bcc}, or prevents the corresponding decision or observation state from satisfying Definitions~\ref{def:z_pf} or~\ref{def:y_cl}, the entire information state $y$ or $z$ should be removed.

%
% 1   考虑Y和Z是耦合的，而Z直接影响系统的可行性，Y只需要保证有通路即可
% 2   首先标记那些不可行的observation state，通过decision state判断
% 2.2 此时infeasible Z可以分为两类：一类是all所有结束状态为deadlock
% 2.3                             一类是存在exit states
% 3   首先删除2.2和2.3
% 4   再根据consistency倒着删，此时一边确定系统是否consistent，一边确定Y是否cognitively live，直到系统CNB
%
%
% 我知道最后的结果就是系统会很破碎，但是：不行的状态就是不行
To synthesize a supervisor that guarantees that the closed-loop system is cognitive non-blocking, we need to prune those states that do not satisfy the properties of potential feasibility, cognitive liveness, or consistency.
%
Although the accessibility of decision macrostates and observation macrostates is coupled, which determines the non-blocking property of the system during its evolution,  any observation state that is a deadlock or livelock within an unobservable reach is clearly infeasible. 
Thus, we start by identifying observation macrostates that are not potentially feasible. 
Specifically, we first detect the observation states corresponding to deadlocks and livelocks, as defined in \eqref{eq:bcc_dl} and \eqref{eq:bcc_ll}, and then prune the associated observation macrostates from the BCC.
%
The removal of infeasible observation macrostates may affect the consistency and cognitive liveness of predecessor decision macrostates. 
For each predecessor decision macrostate $w_y = x_y \in W^Y$ of a pruned observation macrostate, 
if $w_y$ is no longer cognitively live, i.e., \eqref{eq:bcc_y_cl} does not hold for some $q \in x_y$, or it violates Definition~\ref{def:con_bcc} (i), then $w_y$ should also be pruned.
%
After updating decision macrostates, some remaining observation macrostates may become inconsistent. 
We recursively prune inconsistent observation macrostates and their predecessor non-cognitively-live decision macrostates until no further removals are necessary and the system becomes consistent. 
%
We summarize this procedure in Algorithm~\ref{alg:cnb_syn} and illustrate it with Example~\ref{ex:cnb_removal}.
%
% TODO
% 这部分的正确性的concern的担忧主要来自于删除顺序，因为这个问题比原有的IS-based问题在删除过程中会有更强的耦合性，所以得看怎么删
% 最好的办法是过一遍代码检查一下
Note that pruning may leave some macrostates inaccessible from the initial state.
%
Therefore, Line~7 applies the operator $Accessible(\cdot)$ to discard these unreachable remnants and retain only the component accessible from $y_0$.
The resulting BCC $\mathfrak{S}_c$ must also remain macrostate-consistent, meaning that all states associated with each retained decision or observation macrostate are preserved.
This accessible component constitutes a valid cognitive non-blocking supervisor only if it contains at least one cognitively marked decision macrostate in $Y_m$.
Otherwise, the algorithm returns ``No solution''.
%
In the remainder of this article, to simplify the expression, we use $\mathfrak{S}$ instead of $\mathfrak{S}_c$ to denote the consistent cognitive non-blocking BCC output by Algorithm~\ref{alg:cnb_syn}.

\begin{algorithm}[t!]
\caption{Synthesis of Cognitive Non-Blocking BCC}
\label{alg:cnb_syn}
\KwIn{BCC $\mathfrak{S} = (Y, Z, \Sigma, \Gamma, \delta_{YZ}, \delta_{ZZ}, \delta_{ZY}, y_0, Y_m)$\;}
\KwOut{Cognitive non-blocking BCC $\mathfrak{S}_c$\;}

\BlankLine

% \tcp{Step 1: Prune infeasible observation states}
Identify all observation states in $Z_{nl} \subseteq Z$ that are deadlock or livelock using \eqref{eq:bcc_dl} and \eqref{eq:bcc_ll}\;
Prune the observation macrostates $w_z$ corresponding to $(z, q) \in Z_{nl}$\;

\Do{BCC $\mathfrak{S}$ not consistent}{

    Prune the observation states that violate Definition~\ref{def:con_bcc} (i)\;
    Prune the decision states that are not cognitively live or violate Definition~\ref{def:con_bcc} (ii)\;

}

$
\mathfrak{S} = (Y_c, Z_c, \Sigma_c, \Gamma_c, \delta_{YZ,c}, \delta_{ZZ,c}, \delta_{ZY,c}, y_0, Y_{m,c}) \gets Accessible(\mathfrak{S})
$\;
\If{$y_0^Y \notin Y_c$ or $Y_{m,c} = \emptyset$}{
    \Return No solution\;
}
\Return BCC $\mathfrak{S}_c$\;
\end{algorithm}

\begin{example}\label{ex:cnb_removal}
    Let us consider the DFA in Fig.~\ref{fig:dfa_grid_env} (b) and the corresponding BCC in Fig.~\ref{fig:bcc_overall} again.
    We use Algorithm~\ref{alg:cnb_syn} to prune the states that do not satisfy the properties of potential feasibility and cognitive liveness, and then we make this system consistent.
    %
    First, following the result of Example~\ref{ex:y_z_prop}, one can easily check that the blocking observation macrostates $w_{z,2}$, $w_{z,3}$, and $w_{z,6}$ do not satisfy the definition of potential feasibility, and should be pruned.
    We also prune $w_{z,11} = (x_{z,11}, \gamma_{z_11})$ and $w_{z,12} = (x_{z,12}, \gamma_{z_12})$ because their observation states $(x_{11}, q_{19}, \gamma_{11})$ and $(x_{12}, q_{15}, \gamma_{12})$ are deadlock states.
    %
    Next, based on the definition of consistency, we iteratively remove $w_{y,5}$, $w_{z,7}$, and $w_{z,8}$.
    %
    As for decision macrostates $w_{y,1}$ and $w_{y,3}$, one can easily check that there is no feasible sequence that originates from the corresponding decision states and ends at cognitively marked decision macrostate $w_{y,2}$, which violates the requirement of cognitive liveness.
    We need to prune these states, and then we prune $w_{z,4}$ for consistency.
    Finally, based on Definition~\ref{def:con_bcc} (i), all enabled observations must be defined for observation macrostates, so the observation macrostates $w_{z,1}$ should be removed.
    %
    The removed macrostates are marked by red crosses in Fig.~\ref{fig:bcc_overall}, and the remaining structure forms the resulting consistent cognitive non-blocking BCC.
\end{example}


% TODO
% 因为这么删确实会把原系统弄得很破碎
% \begin{remark}
% The BCC obtained by Algorithm~\ref{alg:cnb_syn} may become ``fragmented'' after pruning. This is natural since cognitive non-blocking is stronger than classical non-blocking.
% It requires not only reachability of marked states but also observation-based certification of task completion for those marked states.
% Therefore, some non-blocking states may still be removed. 
% % For supervisor synthesis, we simply restrict attention to the accessible component of the pruned BCC, namely, the part reachable from the initial state $y_0$.
% \end{remark}


%
% TODO
% 其实这图我感觉没啥用但是也就是好看
\begin{figure*}
    \centering
    \includegraphics[width=0.65\linewidth]{Figs/conceptual_proof.pdf}
    \caption{Conceptual illustration of the structural soundness in Theorem~\ref{theo:cnb}. 
    A retained BCC run projects to a feasible DFA run through $P_\Sigma$. For an arbitrary reachable state, a cognitively live decision state admits a continuation to a marked decision state $y_m\subseteq Q_m$, while a potentially feasible observation state admits an observable exit to such a live decision state. 
    Red-crossed components indicate cognitive-blocking or infeasible branches that are excluded from any BCC satisfying Definitions~\ref{def:con_bcc}, \ref{def:z_pf}, and \ref{def:y_cl}.}
    \label{fig:conceptual_proof}
\end{figure*}

\subsection{Synthesized supervisor and correctness}

In this subsection, we show how the cognitive non-blocking BCC $\mathfrak{S}$ induces a supervisor and establishes the correctness of the synthesis algorithm.
%
% Recall that the projection from the BCC language to the sequence of control patterns is
% $
%     P_{\Gamma}:\Sigma_{\mathfrak{S}}^* \rightarrow \Gamma^*,
% $
% and the set of event sequences generated by BCC is denoted as $\mathcal{L}_{\Sigma}(\mathfrak{S})$.
For a BCC run
$
\rho_\mathfrak{S} =
p_0 p_1 \cdots p_i \cdots,
$
$p_i \in P = Y \cup Z$, the corresponding event sequence and control-pattern sequence are
$
    P_\Sigma(tr_\mathfrak{S}(\rho_\mathfrak{S})) \in \Sigma^*,
$
and
$
    P_\Gamma(tr_\mathfrak{S}(\rho_\mathfrak{S}))\in\Gamma^*,
$
respectively.
%
%
Note that a conventional partial-observation supervisor $S: P(\mathcal{L}(G)) \rightarrow \Gamma$ assigns a control pattern $\gamma \in \Gamma$ based on the given observation sequence. 
In the BCC framework, we use $S_{\mathfrak{S}} : P(\mathcal{L}_{\Sigma}(\mathfrak{S})) \rightarrow P_\Gamma(\mathcal{L}(\mathfrak{S}))$ in place of $S$ as a history-based representation of the supervisor induced by the BCC. 
Specifically, $S_{\mathfrak{S}}$ maps each projected event sequence to the sequences of admissible control patterns from the BCC $\mathfrak{S}$, rather than the point-wise control pattern. 
For each observation sequence $\alpha \in P(\mathcal{L}_{\Sigma}(\mathfrak{S}))$, with a slight abuse of notation, define
%
% \begin{flalign}
%     S_{\mathfrak{S}}(\alpha) :=\{&\gamma\in\Gamma\mid \exists \tau \in \mathcal L(\mathfrak{S}): P(P_\Sigma(\tau))=\alpha, \notag \\
%     &\tau\gamma\in\mathcal L(\mathfrak{S})\}
% \end{flalign}
% 
\begin{flalign}\label{eq:bcc_sc:1}
%
% 试过了，这个就是一行放不下
S_{\mathfrak{S}}(\alpha) = \{ &\nu \in \Gamma^* \mid\ \exists \tau \in \mathcal{L}(\mathfrak{S}) : \notag \\
&P(P_\Sigma(\tau)) = \alpha, \nu = P_\Gamma(\tau) \}.
%
% //
% S_{\mathfrak{S}}(\alpha) :=\{&\gamma\in\Gamma\mid \exists \tau \in \mathcal{L}(\mathfrak{S}): P(P_\Sigma(\tau))=\alpha, \notag \\
% &\tau\gamma\in\mathcal{L}(\mathfrak{S})\}.
\end{flalign}
%
%
That is, $S_{\mathfrak{S}}(\alpha)$ collects all sequences of control patterns from those BCC runs whose projected event sequences can generate the same observation $\alpha$.
%
The closed-loop language induced by $S_{\mathfrak{S}}$ is then defined by the event projections of the same BCC runs
\begin{flalign}\label{eq:bcc_sc:2}
\mathcal{L}(S_{\mathfrak{S}}/G)
:=&
\{\, s\in\Sigma^* \mid
\exists \tau \in\mathcal{L}(\mathfrak{S}) :  s=P_\Sigma(\tau),  \notag\\
& \; \; P_\Gamma(\tau) \in S_{\mathfrak{S}}(P(s))
\,\} \notag \\
=&
P_\Sigma(\mathcal{L}(\mathfrak{S}))
=
\mathcal{L}_{\Sigma}(\mathfrak{S}).
\end{flalign}
%
By \eqref{eq:bcc_sc:1}, the condition $P_\Gamma(\tau) \in S_{\mathfrak{S}}(P(s))$ means that the projection onto control pattern for the same BCC run $\tau$ is admissible under the observation $P(s)$, while $s = P_\Sigma(\tau)$ extracts the corresponding event projection.
Thus, \eqref{eq:bcc_sc:1} gives the interpretation of the control sequence for a BCC run, whereas \eqref{eq:bcc_sc:2} extracts its event-sequence interpretation.
%
By Proposition~\ref{prop:bcc_dfa_consistency}, every such projected event sequence is a feasible execution of the DFA $G$.
Hence, we can apply $S_\mathfrak{S}$ as $S$ such that $\mathcal{L}(S_{\mathfrak{S}} / G)$ can be used as an equivalent closed-loop representation of the controlled DES, $\mathcal{L}(S/G)$.
%
Now, we give the following two theorems to show the structural characterization of cognitive non-blockingness. 
%
% Soundness of the algorithm
\begin{theorem}\label{theo:cnb}
Given a DFA $G = (Q, \Sigma, \delta, q_0, Q_m)$ and a corresponding BCC
$
\mathfrak{S} = (Y, Z, \Sigma, \Gamma, \delta_{YZ}, \delta_{ZZ}, \delta_{ZY}, y_0, Y_m),
$
suppose that all decision and observation states of ${\mathfrak{S}}$ are accessible from $y_0$ and satisfy Definitions~\ref{def:con_bcc}, \ref{def:z_pf}, and \ref{def:y_cl}.
The supervisor $S$ induced by ${\mathfrak{S}}$ guarantees that the closed-loop system $S/G$ is cognitive non-blocking.
\end{theorem}
%
\begin{proof}
The intuition of the proof is illustrated in Fig.~\ref{fig:conceptual_proof}. The proof consists of two steps.
First, we use the BCC--DFA consistency to identify the closed-loop behavior $S/G$ with the event projection of the BCC language.
Second, we show that the retained BCC structure is co-accessible, in the cognitive sense, to a marked decision macrostate.

By the definition of the supervisor induced by $\mathfrak{S}$, the closed-loop language generated by $S / G$ coincides with the event projection of the
BCC language, i.e.,
\begin{equation}\label{eq:proof_cnb_lang_equiv}
    \mathcal{L}(S/G) = \mathcal{L}_{\Sigma}(\mathfrak{S})
                     = P_{\Sigma}(\mathcal{L}(\mathfrak{S})).
\end{equation}
Moreover, by Proposition~\ref{prop:bcc_dfa_consistency}, every BCC run projects to a feasible execution of the DFA. Hence, for any $\tau \in \mathcal{L}(\mathfrak{S})$ such that
$
    \delta_{\mathfrak{S}}(y_0, \tau) = p,
$
where $y_0 = \{ q_0 \}$ and $p$ equals to $(x, q) \in Y$ or $(x, q, \gamma) \in Z$,
we have
$
    \delta(q_0, P_\Sigma(\tau)) = q.
$
Therefore, it is sufficient to prove the cognitive non-blocking condition on the BCC and then project the result to the DFA.

Since $\mathfrak{S}$ is accessible, for an arbitrary state $p$ equal to $(x, q)$ or $(x, q, \gamma)$, there exists a sequence $\tau_{q_0, q}$ in $\mathfrak{S}$ such that
$$
p = \delta_{\mathfrak{S}}(y_0, \tau_{q_0, q}),
$$
%
and the corresponding event sequence is denoted by
$P_\Sigma(\tau_{q_0, q}) = s$.
%
We next show that this reachable state admits a cognitively marked continuation. 
By the structural co-accessibility guaranteed in the BCC, there exists a sequence $\kappa_{q, q_m}$ such that
$$
y_m = \delta_{\mathfrak{S}}(p, \kappa_{q, q_m}),
$$
and the corresponding event sequence is denoted by $P_\Sigma(\kappa_{q, q_m}) = t$.
%
If $p$ is a observation state, then Definition~\ref{def:z_pf} guarantees that it can reach a decision state of a successor decision macrostate.
%
For a decision state, \eqref{eq:bcc_y_cl} guarantees that, whenever its macrostate is cognitively live, all corresponding decision states can reach a cognitively marked state through continuations following the same sequence of decision and observation macrostates.
%
Moreover, the bipartite structure preserves observation consistency among decision or observation states with the same first component.
Hence, for any state $p'$ equal to $(x', q')$ or $(x', q', \gamma')$ that shares identical observations with $p$, such that
$$
\begin{aligned}
    p' &= \delta_{\mathfrak{S}}(y_0, \tau_{q_0, q'}) \\
    y_m &= \delta_{\mathfrak{S}}(p, \kappa_{q', q_m}),
\end{aligned}
$$
with
$
P(P_\Sigma(\tau_{q_0, q'})) = P(P_\Sigma(\tau_{q_0, q}))
$
and
$
P(P_\Sigma(\kappa_{q', q_m})) = P(P_\Sigma(\kappa_{q, q_m})).
$
Consequently, for any event sequence 
$s \in \mathcal{L}_{\Sigma}(\mathfrak{S})$, there exists a continuation 
$t \in \Sigma^*$ such that the BCC-level cognitive state estimate induced by the observation $P(st)$ satisfies
%
$$
\begin{aligned}
%\widehat{\bm{\mathcal{E}}}_{\mathfrak{S}}(P(st)) := 
\{ &p \in P \mid 
\exists \tau \in \mathcal{L}(\mathfrak{S}) :  
r = P_{\Sigma}(\tau) \in \Sigma^*\Sigma_o, \\
& \delta_{\mathfrak{S}}(y_0, \tau) = p,\
P(r) = P(st) \}
\subseteq Y_m .
\end{aligned}
$$
%
%Here, $\bm r$ is a BCC sequence, while $r=P_\Sigma(\bm r)$ is its projected DFA event sequence. 
Hence, the above inclusion states that all BCC states consistent with the same observation $P(st)$ are cognitively marked.
%
By \eqref{eq:proof_cnb_lang_equiv} and Proposition~\ref{prop:bcc_dfa_consistency}, the above BCC-level inclusion can be projected back to the corresponding closed-loop DFA $S_{\mathfrak{S}} / G$. Since every state in $Y_m$ is of the form $(x_{y,m}, q_{y,m})$ with $x_{y,m} \subseteq Q_m$, its DFA component satisfies $q_{y,m} \in Q_m$. 
Therefore, for any observation sequence $\alpha \in P(\mathcal{L}(S_{\mathfrak{S}}/G))$, there exists a continuation $\omega \in \Sigma_o^*$ such that the corresponding cognitive state estimate satisfies 
$
\hat{\mathcal{E}}_{S_{\mathfrak{S}}}(\alpha\omega) \subseteq Q_m .
$
This completes the proof.
\end{proof}

% Soundness: 只要系统证明了某件事，那么这件事在语义上确实成立。
% Completeness: 只要某件事实际上成立，系统就一定能把它证明/表示出来。
\begin{theorem}\label{theo:cnb_complete}
Given a DFA $G = (Q, \Sigma, \delta, q_0, Q_m)$, suppose that the closed-loop system $S / G$ satisfies the notion of cognitive non-blocking.
Then, the corresponding BCC
$
\mathfrak{S} = (Y, Z, \Sigma, \Gamma, \delta_{YZ}, \delta_{ZZ}, \delta_{ZY}, y_0, Y_m),
$
constructed for $S / G$ satisfies Definitions~\ref{def:con_bcc},
\ref{def:z_pf}, and~\ref{def:y_cl}.
\end{theorem}

\begin{proof}
By Definition~\ref{def:cnb}, for any $s \in \mathcal{L}(S / G)$ with $q = \delta(q_0, s)$, there exists a continuation $t$ such that $st \in (\Sigma^*\Sigma_o) \cap \mathcal{L}(S / G)$ and, for any $t' \in (\Sigma^*\Sigma_o) \cap \mathcal{L}(S / G)$,
\begin{equation}\label{eq:theo_comp:1}
    P(t') = P(st) \Rightarrow \delta(q_0, t') \in Q_m.
\end{equation}
%
We now examine how the reachable state $q$ appears in the corresponding evolution of marcostates. 
If $q = q_0$, then $q$ belongs to the initial decision macrostate $w_{y,0} = x_{y,0}$, and the following decision macrostate argument applies directly.

First, suppose that the incoming transition of $q$ is observable and can be written as $q' \xrightarrow[]{\sigma_o} q$. 
Then there exists an observation macrostate $w_z' = (x', \gamma)$ with $q' \in x'$ such that
$$
NX_{\sigma_o}(x') = x_y, \quad q \in x_y,
$$
i.e., $q$ belongs to a decision macrostate $w_y = x_y$.
%
By \eqref{eq:theo_comp:1}, after the continuation $t$, all closed-loop executions observation-equivalent to $st$ end in marked states. 
Since observation-equivalent executions are collected into the same successor decision macrostate, the decision macrostate $w_{y_m}' = x_{y,m}'$ reached after $P(st)$ satisfies
$
x_{y,m}' \subseteq Q_m.
$
Thus, $w_{y_m}'$ is cognitively marked, and $w_y$ is cognitively live.
Hence, Definition~\ref{def:y_cl} holds.

Second, suppose that the incoming transition of $q$ is unobservable and can be written as $q' \xrightarrow[]{\sigma_{uo}} q$. Then there exists a decision state $y'$ and a control pattern $\gamma$ such that
$$
x_z = UR_{\gamma}(y') \wedge q_z \in x_z,
$$
i.e., $q_z$ belongs to the observation macrostate $w_z = (x_z, \gamma_z)$ and to the observation state $(x_z, q, \gamma_z)$.
%
If $w_z$ contained a deadlock observation state, then the corresponding closed-loop sequence would admit no continuation to a later observation certifying \eqref{eq:theo_comp:1}. 
If $w_z$ contained a livelock consisting only of the same observation state, then the system could stay within the unobservable reach of $w_z$ forever without reaching a successor decision macrostate, again contradicting \eqref{eq:theo_comp:1}. 
Therefore, no reachable observation macrostate in $\mathfrak{S}$ contains such deadlock or livelock observation states, and Definition~\ref{def:z_pf} holds.

It remains to verify consistency. 
Since ${\mathfrak{S}}$ is induced by the closed-loop behavior of a partial-observation supervisor, every reachable decision macrostate has at least one control decision selected by $S$. 
Moreover, for each reachable observation macrostate, every feasible observable event enabled under its control pattern must be represented by a successor decision macrostate.
Otherwise, some possible closed-loop evolution under the same observation would be missing from $\mathfrak{S}$. 
Hence, Definition~\ref{def:con_bcc} also holds.
%
Consequently, the BCC constructed for the cognitive non-blocking closed-loop system satisfies Definitions~\ref{def:con_bcc}, \ref{def:z_pf}, and~\ref{def:y_cl}.
\end{proof}

Intuitively, a closed-loop system is cognitively non-blocking if and only if its behavior can be represented by an accessible BCC whose behavior satisfies the requirements of consistency, potential feasibility, and cognitive liveness.
%
We also give the following corollary for the correctness of Algorithm~\ref{alg:cnb_syn}.
The algorithm returns a nonempty BCC exactly when the cognitive non-blocking synthesis problem is solvable. 
If a nonempty BCC remains after pruning, it directly induces a cognitive non-blocking supervisor. 
Conversely, if such a supervisor exists, then its closed-loop behavior admits a BCC representation satisfying the same structural requirements, and hence the pruning procedure cannot eliminate all states.

\begin{corollary}\label{cor:alg_cnb_sound}
Algorithm~\ref{alg:cnb_syn} returns a nonempty BCC rather than ``No solution'' if and only if there exists a partial-observation supervisor $S$ such that the closed-loop system $S/G$ is cognitively non-blocking.
%
Moreover, if Algorithm~\ref{alg:cnb_syn} returns a nonempty BCC $\mathfrak{S}$, then the supervisor $S_{\mathfrak{S}}$ induced by $\mathfrak{S}$ guarantees that $S_{\mathfrak{S}}/G$ is cognitive non-blocking.
\end{corollary}

% \begin{proof}
% We first prove the ``only if'' part of this corollary.
% Suppose that Algorithm~\ref{alg:cnb_syn} returns a nonempty BCC $\mathfrak{S}$. 
% By construction, the algorithm removes all observation macrostates that are not potentially feasible, all decision macrostates that are not cognitively live, and all decision or observation macrostates that violate the consistency requirement. It then extracts the accessible component from the initial decision state $y_0$. 
% Hence, $\mathfrak{S}$ is accessible from $y_0$ and satisfies Definitions~\ref{def:con_bcc}, \ref{def:z_pf}, and~\ref{def:y_cl}. 
% By Theorem~\ref{theo:cnb}, the supervisor $S_{\mathfrak{S}}$ induced by $\mathfrak{S}$ guarantees that $S_{\mathfrak{S}}/G$ is cognitive non-blocking. 
% Therefore, there exists a partial-observation supervisor enforcing cognitive non-blocking.

% We next prove the ``if'' part. 
% Suppose that there exists a partial-observation supervisor $S$ such that the closed-loop system $S/G$ is cognitively non-blocking. 
% Let ${\mathfrak{S}}'$ be the corresponding BCC constructed for $S / G$. 
% By Theorem~\ref{theo:cnb_complete}, ${\mathfrak{S}}'$ satisfies Definitions~\ref{def:con_bcc}, \ref{def:z_pf}, and~\ref{def:y_cl}.
% %
% We show that Algorithm~\ref{alg:cnb_syn} cannot remove all states of ${\mathfrak{S}}'$. 
% Consider the BCC ${\mathfrak{S}}'$ obtained after each pruning iteration, and maintain the invariant component of ${\mathfrak{S}}'$ that all states of ${\mathfrak{S}}'$ are still contained in the current BCC. 
% One can easily check that this invariant component holds initially.
% %
% Suppose that the invariant component holds before one pruning iteration. 
% Since the observation macrostates within ${\mathfrak{S}}'$ satisfy Definition~\ref{def:z_pf}, none of its observation states can be removed as an infeasible observation state. 
% Considering the decision macrostates ${\mathfrak{S}}'$ satisfy Definition~\ref{def:y_cl}, none of its decision states can be removed as a non-cognitively-live decision state. 
% Meanwhile, since ${\mathfrak{S}}'$ satisfies Definition~\ref{def:con_bcc}, none of its decision or observation states can be removed due to a consistency violation. 
% Thus, the invariant component is preserved after this iteration.

% Therefore, after the pruning procedure terminates, the remaining BCC still contains the accessible BCC representation ${\mathfrak{S}}'$. 
% In particular, the accessible component from $y_0$ is nonempty and contains a cognitively marked decision state. 
% Finally, Algorithm~\ref{alg:cnb_syn} cannot return ``No solution''.
% \end{proof}
%

\begin{figure}[ht!]
    \centering
    \includegraphics[width=0.55\linewidth]{Figs/controlled_fsa_overall.pdf}
    \caption{Controlled DFA $S / G$ from DFA $G$ in Fig.~\ref{fig:dfa_grid_env} and BCC $\mathfrak{S}$ in Fig.~\ref{fig:bcc_overall}.}    % TODO 这个符号对不对
    \label{fig:ctrl_DFA}
\end{figure}

We give Example~\ref{ex:DFA_cnb} to illustrate that the closed-loop system satisfies the notion of cognitive non-blocking.
Meanwhile, we discuss the complexity of Algorithm~\ref{alg:cnb_syn}.

\begin{example}\label{ex:DFA_cnb}
    Let us consider the DFA illustrated in Fig.~\ref{fig:dfa_grid_env} (b) again, and the corresponding BCC illustrated in Fig.~\ref{fig:bcc_overall}.
    %
    After pruning the inconsistent states, the corresponding closed-loop system $S / G$ is illustrated in Fig.~\ref{fig:ctrl_DFA}.
    %
    In the resulting closed-loop system, the evolutions from the initial state $q_0$ are $q_0 \xrightarrow{u_1} q_1 \xrightarrow{o_2} q_5$ and $q_0 \xrightarrow{u_2} q_2 \xrightarrow{o_2} q_6$, and the corresponding observation is $o_2$.
    In the cognitive sense, after observing $o_2$, the supervisor can certify that the current DFA state belongs to the marked set
    \begin{equation}\label{eq:ex_cnb_reach_marked}
        \hat{\mathcal{E}}_S(o_2) = \{ q_5, q_6 \} \subseteq Q_m,
    \end{equation}
    although the exact state may still be unknown.
    %
    Moreover, from any reachable state after $q_5$ or $q_6$, the controlled system can evolve through a finite admissible continuation and return to $q_0$.
    The particular observable labels along this return segment are not essential for the cognitive non-blocking argument.
    After returning to $q_0$, the continuation for states of \eqref{eq:ex_cnb_reach_marked} can be applied again to reach the cognitively marked information state $\{ q_5, q_6 \}$.
    %
    Hence, for any sequence in $G$ corresponding to a cognitively live path in the BCC $\mathfrak{S}$, there exists a continuation that reaches a cognitively marked state.
    Once the supervisor $S$ is induced from the BCC $\mathfrak{S}$ as described in Section~\ref{subsec:supervisor_induction}, the closed-loop system $S/G$ inherits this property and satisfies cognitive non-blocking.
    
    Let us consider the motivating example again in Section~\ref{sec:intro}.
    The robotic system cannot certify its exact current state, but its state estimate can be updated after completing certain motion primitives.
    When it is at the initial state $q_0$, it should reject the motions corresponding to $o_3$, which may lead to $q_3$ or $q_4$.
    In this sense, it can then arrive at $q_5$ or $q_6$; moreover, when it observes the information corresponding to event $o_2$, it knows that it is in a safe location, so the transmission of critical information can be started.
\end{example}


% \begin{proposition}
% Consider a BCC $\mathfrak{S}$ constructed over a DFA $G$ with the number of $|Q|$ states and $|\Sigma|$ events. 
% The worst-case time complexity of Algorithm~\ref{alg:cnb_syn} is
% $
%     O(|Q|\cdot |\Sigma|\cdot 2^{2|Q|+2|\Sigma|}).
% $
% \end{proposition}
%
% \begin{proof}
% Algorithm~\ref{alg:cnb_syn} proceeds by pruning decision and observation macrostates of the BCC.
% In Lines 1--2, deadlock and livelock observation state are identified by a BFS
% search, which can be done in time
% $
% O(|Y + Z|+|\delta_{\mathfrak{S}}|)
% $.
% %
% In Lines 3--6, each iteration prunes decision macrostates that are not cognitively live
% or observation macrostates that violate consistency. Each iteration can be implemented by one BFS over the remaining BCC, with cost
% $
% O(|Y + Z|+|\delta_{\mathfrak{S}}|)
% $.
% Since each nontrivial iteration prunes at least one decision macrostate or observation macrostate, the
% number of iterations is at most $|W^Y|+|W^Z|$. Therefore, the total cost is
% $
%     O((|W^Y|+|W^Z|)(|Y + Z|+|\delta_{\mathfrak{S}}|)).
% $
% Line 7 performs one additional BFS, which is dominated by the above term.
% Hence, the overall complexity is
% $
%     O((|W^Y|+|W^Z|)(|Y + Z|+|\delta_{\mathfrak{S}}|)).
% $
% %
% %By Proposition~\ref{prop:state_bcc}, we have
% By \eqref{eq:prop:bcc:1} and \eqref{eq:prop:bcc:2}, we have
% $
%     |Y + Z|
%     \le
%     |Q|(2^{|Q|}+2^{|Q|+|\Sigma|})
%     =
%     O(|Q|\cdot 2^{|Q|+|\Sigma|}),
% $
% and
% $$
%     |\delta_{\mathfrak{S}}|
%     \le
%     |Q|2^{|Q|+|\Sigma|}(1+|\Sigma|)
%     =
%     O(|Q|\cdot |\Sigma|\cdot 2^{|Q|+|\Sigma|}).
% $$
% Therefore, based on
% $
% |W^Y|+|W^Z| = O(2^{|Q|+|\Sigma|}),
% $
% we have
% $$
% \begin{aligned}
%     O((|W^Y|+|W^Z|)(|Y + Z|+|\delta_{\mathfrak{S}}|))
%     &=
%     O(
%     2^{|Q|+|\Sigma|}
%     \cdot
%     |Q|\cdot |\Sigma|\cdot 2^{|Q|+|\Sigma|}
%     ) \\
%     &=
%     O(|Q|\cdot |\Sigma|\cdot 2^{2|Q|+2|\Sigma|}).
% \end{aligned}
% $$
% This completes the proof.
% \end{proof}
%

The pruning-based synthesis procedure in Algorithm~\ref{alg:cnb_syn} is exponential in both the state space and the event alphabet of the underlying DFA.
Its worst-case computational complexity is given by
$O(|Q|\cdot |\Sigma|\cdot 2^{2|Q|+2|\Sigma|})$, where $|Q|$ is the number of states, and $|_\Sigma|$ is the number of events.

As for the maximally permissiveness of the synthesized supervisor, the comparison of control patterns should be made at the same decision instant. Let
$
    \iota = w_{y,0} w_{z,0} w_{y,1} w_{z,1} \cdots w_{y,n} w_{z,n} \cdots
$
be an information-state sequence of BCC $\mathfrak{S}$ generated by the synthesized supervisor \eqref{eq:bcc_sc:1}, where $\gamma_i$ is the control pattern made at decision macrostate $w_{y,i}$.
The maximally permissiveness requires that for any other $\gamma_i' \in \mathcal{H}(w_{y,i})$, we have $\gamma_i \not \subset \gamma_i'$.
%
In this sense, the maximally permissive partial-observation supervisor can be defined as
\begin{flalign}\label{eq:mps_cnb}
    S(\alpha) = \{ &\nu = \gamma_0 \gamma_1 \cdots \gamma_n \cdots \in \Gamma^* \mid \exists \rho_{\mathfrak{S}} \in \mathcal{R}_{\mathfrak{S}}, 
    \tau = tr_{\mathfrak{S}}(\rho_{\mathfrak{S}}) :  \notag \\
    &\alpha = P(P_\Sigma(\tau)), \nu = P_\Gamma(\tau), \iota = w_{y,0} w_{z,0} w_{y,1} w_{z,1} \cdots w_{y,n}   \notag \\
    & w_{z,n} \dots = R_X(\rho_{\mathfrak{S}}),
    \forall \gamma_i' \in \mathcal{H}(w_{y,i}), \gamma_i \not \subseteq \gamma_i'; \},
\end{flalign}
where $i \in \mathbb{N}$.

\begin{corollary}
Problem~\ref{prob:cnb} can be solved within finite steps by Algorithm~\ref{alg:cnb_syn}, and the maximally permissive supervisor can be obtained from \eqref{eq:mps_cnb}.
\end{corollary}


\section{Special case: marked tasks within the unobservable reaches}
\label{sec:special_z_marked}

%
% 这个系统要处理的话
% 其实方法是一样的
% 1 对于marked observation state，需要当前Z-状态中的所有状态能够到达这个状态，或者退到下一个Y状态
% 2 Y状态不变
% 3 系统要consistent
% 

In this section, we discuss a special case of the proposed method, namely, systems in which marked tasks may be completed within unobservable reaches.
In the previous sections, cognitive non-blockingness is formulated based on the certification of marked states after observation updates.
This formulation is consistent with the physical interpretation that task completion should be certified from the observation.
However, from a mathematical perspective, this requirement can be restrictive when a marked state is reached within an unobservable evolution before the next observable event occurs.
%
To cover this situation, we extend the notion of cognitive non-blockingness by allowing task completion to be certified even if the marked state is visited within the preceding unobservable reach, and the formal definition is given as follows
\begin{equation}
    \forall s \in \mathcal{L}(S/G), \exists t \in \Sigma^* : \mathcal{E}_S(st) \subseteq Q_m.
\end{equation}

Under this formulation, the cognitively marked states are naturally defined over decision states, since decision states represent information states immediately after observation updates.
However, in some systems, a marked state can be reached inside a observation state before the next observable event occurs.
%
The main result of this article already handles the case where task completion is certified at marked decision states. Therefore, we only briefly discuss the case of marked observatino states here.

To handle this case, we extend the target set of the BCC by including marked observation states, which are formally defined as
$
    Z_{m} = \{ (x_{z,m}, q_m, \gamma) \in Z \mid x_{z,m} \subseteq Q_m \}.
$
Meanwhile, the set of exit observation states can be defined as
$$
    Z_{e}
    =
    \{ (x_{z,e}, q_e, \gamma) \in Z \mid
    \exists \sigma_o \in (\Gamma(z) \cup \Sigma_{uc}) \cap \Sigma_o : \delta(q_e,\sigma_o)! \}.
$$
%
Given a observation macrostate $w_{z,m} = (x_{z,m}, \gamma_{z,m})$, its corresponding set of observation state is $Z_{x_{z,m}, \gamma_{z,m}}$, which shares information state $x_{z,m}$ and control pattern $\gamma_{z,m}$. 
We say that $w_{z,m}$ is a \emph{marked observation macrostate} if (i) for every observation state $z \in Z_{x_{z,m}, \gamma_{z,m}}$, there exists a marked observation state $z_m \in Z_{x_{z,m}, \gamma_{z,m}} \cap Z_{m}$ such that
$$
    z \leadsto_{Z_{x_{z,m}, \gamma_{z,m}}} z_m.
$$
(ii) for this marked observation state $z_m$, every path originating from $z_m$ and inside $Z_{x_{z,m}, \gamma_{z,m}}$ must end with an exit observation state $z_e \in Z_{e}$, that is, for every
$$
    \rho_{\mathfrak{S}, i} = z_{i,0} \ z_{i,1} \ \cdots \ z_{i,j} \ \cdots \ z_{i,\ell},
$$
$z_{i,0} = z_m, \, z_{i,\ell} = q_e$, we have
$
    z_{i,j} \leadsto_{Z_{x_{z,m}, \gamma_{z,m}}} z_{k,e},
$
with $i = 0, \ldots , \ell - 1$, $j,k \in \mathbb{N}$.
Hence, after reaching the marked observation state $z_m$, all subsequent internal evolution inside $Z_{x_{z,m}, \gamma_{z,m}}$ can eventually reach an exit observation state $z_{k,e}$.
%
Thus, the observed exit can certify task completion within the unobservable reach.

Then, the target set for cognitive-liveness is extended from
$Y_m$ to
\begin{equation}
    P_m^{+}
    =
    Y_m \cup \bigcup_{w_{z,m} \in W^Z} Z_{x_{z,m}, \gamma_{z,m}},
    \label{eq:extended_target_set}
\end{equation}
where $w_{z,m} = (x_{z,m}, \gamma_{z,m})$.
Accordingly, the definition of cognitively live decision macrostates $w_{y, l} = x_{y,l}$ in \eqref{eq:bcc_y_cl} is extended as
\begin{equation}
\label{eq:bcc_y_cl:2}
    \forall q_l \in x_{y,l},\;
    \exists p \in P_m^{+} :
    p = \delta_{\mathfrak{S}}((x_{y,l}, q_l),tr_{\mathfrak{S}}(\rho_{\mathfrak{S}}')),
\end{equation}
where the same BCC sequence $tr_{\mathfrak{S}}(\rho_{\mathfrak{S}}')$ is through the set of potentially feasible observation macrostates, and the new cognitively live macrostates, as in \eqref{eq:bcc_y_cl}.
Thus, a cognitively live decision state is allowed to reach either a cognitively marked decision state or a marked observation state with an internal path to the next observable exit.
%
Based on these modifications, Algorithm~\ref{alg:cnb_syn} can still be applied to obtain a consistent BCC. 
The resulting BCC then induces a cognitive non-blocking supervisor for systems in which task completion can be certified either at marked decision states or within marked observation states.
In this way, Problem~\ref{prob:cnb} can be handled by replacing \eqref{eq:est_o} with \eqref{eq:esti:1}.

It should be noted that completion of a marked observation state is certified only after the subsequent observable exit is observed. 
This is sufficient for the mathematical definition, since all compatible BCC paths must have visited a marked observation state before the exit. 
However, in many practical applications, the unobservable tail inside the observation state is not physically observed, and thus the exact occurrence of task completion may not be directly interpretable from the subsequent observation alone.

\section{Conclusion}\label{sec:con}

In this article, we proposed a new non-blocking property, called cognitive non-blocking, under which the observer can certify, from partial observations, that the system has reached a marked state.
Based on this notion, we developed a supervisor synthesis algorithm that constructs a controlled sublanguage guaranteed to be cognitive non-blocking. 
We also proved that the proposed synthesis algorithm terminates in finite steps. 
%
Future work will investigate how the notion of cognitive non-blocking can be extended to weighted and interval-weighted systems, where task completion may depend not only on reachability but also on accumulated
costs or weights.

% \section*{Acknowledgements}
% The authors thank the anonymous reviewers for their helpful comments.

% A one-file author--year bibliography. No .bib file is needed for this sample.
\bibliographystyle{elsarticle-harv}
\bibliography{ref_full}

\end{document}
